\documentclass[11pt,a4paper]{article}
\usepackage{amsmath,amssymb,mathtools,bm}
\usepackage[final]{microtype}
\usepackage{xurl}
\usepackage{geometry}
\geometry{margin=1in}
\urlstyle{tt}

\title{First-Principles Toroidal Shell Identities\\
and Radial/Projected Closures for Full Induction}
\author{}
\date{}

\begin{document}
\maketitle

\section*{Purpose}

This note rewrites the toroidal shell identities underlying the full-induction
\texttt{dt\_alpha} path from first principles and then compares them to the
current implementation.

The point of the note is to keep three things separate:
\begin{enumerate}
\item exact identities that follow from the stated assumptions;
\item additional model choices used to obtain a closed shell problem;
\item numerical choices used by the implementation.
\end{enumerate}

If a block in the code is not derivable from the assumptions stated here, it
should not be described as ``first-principles'' without an additional derivation.

\paragraph{Scope.}
This is a note about the \emph{local toroidal/FAC shell backbone}. It is not, by
itself, a complete note for the full self-consistent induction problem. A full
closure note must also specify how the upper-region induced response enters
through the driving/response operator.

\paragraph{Two operator languages.}
The note now uses two closely related, but logically distinct, operator
languages:
\begin{enumerate}
\item an \emph{electric/connector} language, organized around
\[
q=\partial_r U-E_r,
\qquad
\mathcal D_q,
\qquad
\Gamma_{\rm known},
\]
which is the natural language when the primitive shell input is the
known tangential electric trace \(\mathbf E_{S,\rm shell}\);
\item a \emph{magnetic reduced-shell} language, organized around
\[
\Pi=\mathcal P_I[\chi],
\qquad
K_I \chi = -R_I \widetilde G,
\]
which is the natural language when the shell toroidal update \(\chi\) itself
is the unknown and the upper response is written as the non-potential
tangential magnetic correction after harmonic-baseline subtraction.
\end{enumerate}
These are not competing closures. The first is the constructive interface that
builds \(q\) from forcing data; the second is the compressed shell statement
once \(\chi\) is the unknown.

\paragraph{Notation bridge to the companion document.}
The companion reduced-shell document uses the symbol \(q\) for the shell
toroidal update. In the present note, that same shell quantity is written
\[
\chi=\partial_t\psi^+ = \frac{1}{R_I}\Pi_0 q,
\]
while the symbol \(q\) is reserved for the electric connector trace
\[
q:=\partial_r U-E_r.
\]
Here \(U\) is the \emph{physical} curl-free shell electric scalar, not the raw
Helmholtz coefficient \(\Phi\). The code represents the tangential electric
field using the repo Helmholtz sign convention
\[
\mathbf E_S
=
s_{\rm cf}\,\nabla_S\Phi
s_{\rm df}\,\hat r\times \nabla_S W,
\qquad
(s_{\rm cf},s_{\rm df})=(-1,-1)
\]
in the current implementation. Throughout this note we therefore set
\[
U := s_{\rm cf}\,\Phi,
\qquad
V := W,
\]
so that
\[
\mathbf E_S = \nabla_S U - \hat r\times \nabla_S V.
\]
With the active repo sign choice \(s_{\rm cf}=-1\), this means
\[
U=-\Phi,
\]
which matches the legacy electrostatic sign convention. The connector scalar
\(q=\partial_r U-E_r\) should be read in that physical-\(U\) convention.
So whenever the companion document writes \(\Pi=\mathcal P_I[q]\), the
corresponding statement here is
\[
\Pi=\mathcal P_I[\chi].
\]

\paragraph{Reading guide.}
The note is organized in three layers.
\begin{enumerate}
\item Sections 2--8 derive the exact shell identities and explain the
relationship between the radial and tangential closures.
\item Sections 9--10 formulate the canonical radial-shell closure, the missing
scalar, and the abstract shell/co-energy reduction.
\item The appendices explain how the present runtime realizes those objects,
what parts are exact for the adopted model, and what parts are still reduced
or explicitly modeled.
\end{enumerate}
So the early tangential analysis should be read mainly as a sign and structure
audit of the older projected-shell route, while the later sections give the
recommended interpretation of the current radial-shell path.
Readers interested only in the canonical radial-shell closure can skim the
projected-tangential detour in Sections 6--8 and go directly from the exact
shell backbone to Section 9.

\section{Assumptions}

We assume only:
\begin{enumerate}
\item linearized quasi-static Maxwell equations;
\item a prescribed static background field $\mathbf B_0$;
\item the strong field-aligned current ansatz
\[
\mathbf j = \alpha \mathbf B_0;
\]
\item current continuity
\[
\nabla\cdot \mathbf j = 0;
\]
\item a thin ionospheric shell at $r=R_I$;
\item a toroidal magnetic shell scalar $\psi$ for the induction update.
\end{enumerate}

No particular upper-region model is assumed yet. Any such model enters only
through the driving electric field or a separate response operator.

\section{Geometry and conventions}

Let $S_I=\{r=R_I\}$ be the ionospheric shell. For a scalar $f$,
\[
\nabla_S f = \frac{1}{r}\nabla_\Omega f,
\qquad
\Delta_S f = \frac{1}{r^2}\Delta_\Omega f.
\]

For a tangential vector $\mathbf a=(a_\theta,a_\phi)$,
\[
\hat r\cdot (\nabla\times \mathbf a)
=
\frac{1}{r}
\left(
\partial_\theta a_\phi + \cot\theta\,a_\phi - \partial_\phi^\ast a_\theta
\right),
\qquad
\partial_\phi^\ast := \frac{1}{\sin\theta}\partial_\phi.
\]

The toroidal magnetic scalar uses the convention
\begin{equation}
\mathbf b_T = \nabla\times (r\psi\,\hat r)
=
-\,\hat r\times \nabla_S(r\psi).
\label{eq:b_tor}
\end{equation}

This is the same sign convention used by the code for the imposed toroidal
scalar and for the toroidal magnetic induction update.

\section{Strong FAC consequences}

Define
\[
\eta := \partial_t \alpha.
\]
Since $\nabla\cdot \mathbf B_0=0$ and $\nabla\cdot(\alpha \mathbf B_0)=0$,
\[
\mathbf B_0\cdot \nabla \alpha = 0,
\qquad
\mathbf B_0\cdot \nabla \eta = 0.
\]
Thus both $\alpha$ and $\eta$ are field-line invariants.

Writing
\[
\mathbf B_0 = B_{0r}\hat r + \mathbf B_{0S},
\qquad
\mathbf B_{0S}=(B_{0\theta},B_{0\phi}),
\]
the shell current is
\begin{equation}
j_r = \alpha B_{0r},
\qquad
\mathbf j_S = \alpha \mathbf B_{0S},
\label{eq:jr_js_alpha}
\end{equation}
and therefore
\begin{equation}
\partial_t j_r = B_{0r}\eta,
\qquad
\partial_t \mathbf j_S = \eta \mathbf B_{0S}.
\label{eq:dtjr_dtjs_alpha}
\end{equation}

\paragraph{Notation reminder.}
Here \(\mathbf j_S\) denotes the tangential \emph{volume-current trace} just
off the shell. Later in the note, when the ionosphere is treated as a thin
current sheet, \(\mathbf K_S\) will denote the corresponding
height-integrated \emph{sheet current}. The two play different roles, and the
distinction matters in the continuity-connector sections.

\section{Exact toroidal shell identities}

\subsection*{Radial current}

From \eqref{eq:b_tor} and Amp\`ere's law,
\begin{equation}
j_r = -\,\frac{1}{\mu_0 R_I}\,\Delta_\Omega \psi.
\label{eq:jr_psi}
\end{equation}
Hence
\begin{equation}
\partial_t j_r
=
-\,\frac{1}{\mu_0 R_I}\,\Delta_\Omega \partial_t \psi.
\label{eq:dtjr_dtpsi}
\end{equation}
Combining \eqref{eq:dtjr_dtjs_alpha} and \eqref{eq:dtjr_dtpsi} gives
\begin{equation}
\Delta_\Omega \partial_t \psi
=
-\,\mu_0 R_I\, B_{0r}\eta.
\label{eq:dtpsi_from_eta_core}
\end{equation}

On the mean-zero subspace this can be written
\begin{equation}
\partial_t \psi
=
-\,\mu_0 R_I\, \Delta_\Omega^{-1}\Pi_0(B_{0r}\eta).
\label{eq:dtpsi_from_eta}
\end{equation}

This is exact. It shows that $\partial_t \psi$ is not an independent unknown
once $\eta$ is known.
The projector \(\Pi_0\) is not only a gauge choice: it also enforces the
mean-zero shell solvability condition needed for the inversion of
\(\Delta_\Omega\) on the sphere.

\subsection*{Tangential current}

The tangential part of $\nabla\times \mathbf b_T$ gives
\begin{equation}
\mathbf j_S
=
+\,\frac{1}{\mu_0}\,\nabla_S\!\bigl(\partial_r(r\psi)\bigr).
\label{eq:js_psi}
\end{equation}
Therefore
\begin{equation}
\mu_0\,\partial_t \mathbf j_S
=
+\,\nabla_S\!\bigl(\partial_r(r\,\partial_t\psi)\bigr)
=
+\,\nabla_S\!\bigl(\partial_t\psi + R_I \partial_r\partial_t\psi\bigr).
\label{eq:dtjs_dtpsi}
\end{equation}

Combining \eqref{eq:dtjr_dtjs_alpha} and \eqref{eq:dtjs_dtpsi},
\begin{equation}
\mu_0\,\eta\,\mathbf B_{0S}
=
+\,\nabla_S\!\bigl(\partial_t\psi + R_I \partial_r\partial_t\psi\bigr).
\label{eq:tangential_vector_balance}
\end{equation}

Projecting \eqref{eq:tangential_vector_balance} onto $\mathbf B_{0S}$ gives
\begin{equation}
\mu_0 |\mathbf B_{0S}|^2 \eta
=
\mathbf B_{0S}\cdot \nabla_S\!\bigl(\partial_t\psi + R_I \partial_r\partial_t\psi\bigr).
\label{eq:tangential_projected_balance_no_driver}
\end{equation}

\paragraph{Important point.}
Equation \eqref{eq:tangential_projected_balance_no_driver} is also an exact
first-principles shell identity. It is \emph{not} the same equation as the
radial shell balance, but it is not arbitrary either. It comes from the
tangential part of the same Maxwell system together with the same strong FAC
ansatz.

\paragraph{But it is weaker than the full tangential equation.}
The vector equation \eqref{eq:tangential_vector_balance} contains two
independent tangential components. Its projection
\eqref{eq:tangential_projected_balance_no_driver} keeps only the component
parallel to $\mathbf B_{0S}$. So the projected scalar equation is an exact
consequence of the full tangential equation, but not conversely.

Therefore using $|\mathbf B_{0S}|^2$ in a toroidal shell equation is not
``replacing $B_{0r}$ by $|\mathbf B_{0S}|$''. Rather:
\begin{itemize}
\item $B_{0r}$ enters the exact $j_r \leftrightarrow \psi$ relation
\eqref{eq:dtpsi_from_eta_core};
\item $|\mathbf B_{0S}|^2$ enters the exact tangentially projected shell
balance \eqref{eq:tangential_projected_balance_no_driver}.
\end{itemize}

These equations become weak in complementary regions:
\begin{itemize}
\item the radial equation is weak where $B_{0r}$ is small;
\item the tangential projected equation is weak where $\mathbf B_{0S}$ is small.
\end{itemize}
Neither one is globally superior by principle alone.

\section{Exact radial closure identities}

The code uses the one-sided radial derivative of $\partial_t j_r$ and
$\partial_t\psi$. These are derivable from the same assumptions.

From field-line invariance,
\[
B_{0r}\partial_r \eta + \mathbf B_{0S}\cdot \nabla_S \eta = 0.
\]
Hence
\begin{equation}
\partial_r(\partial_t j_r)
=
\partial_r(B_{0r}\eta)
=
(\partial_r B_{0r})\eta - \mathbf B_{0S}\cdot \nabla_S \eta.
\label{eq:dr_dtjr}
\end{equation}
Equivalently, in unit-sphere angular derivatives,
\begin{equation}
\partial_r(\partial_t j_r)
=
(\partial_r B_{0r})\eta
-\frac{1}{R_I}\,\mathbf B_{0S}\cdot \nabla_\Omega \eta.
\label{eq:dr_dtjr_unit_sphere}
\end{equation}

Now use the shell relation
\[
\partial_t \psi = -\mu_0 r\,\Delta_\Omega^{-1}(\partial_t j_r)
\]
at general shell radius $r$. Differentiating in $r$ and evaluating at $r=R_I$
gives
\begin{equation}
\partial_r\partial_t \psi
=
\frac{1}{R_I}\,\partial_t\psi
+
T_{j_r\to\psi}\,\partial_r(\partial_t j_r),
\label{eq:dr_dtpsi}
\end{equation}
where
\begin{equation}
T_{j_r\to\psi}
:=
-\,\mu_0 R_I\,\Delta_\Omega^{-1}
\label{eq:T_jr_to_psi}
\end{equation}
on the mean-zero subspace.

Equations \eqref{eq:dtpsi_from_eta}, \eqref{eq:dr_dtjr_unit_sphere}, and
\eqref{eq:dr_dtpsi} are exactly the identities used by the code in the
\texttt{dt\_alpha -> dt\_psi} rewrite.

\section{Driver/response split}

\paragraph{Optional detour.}
This section introduces both the canonical radial-shell closure and the exact
projected tangential scalar closure. Readers whose only goal is the radial-shell
path may read the radial-shell part first and treat the projected-tangential
part, together with the sign audit in the next sections, as supporting
background for the older \path{tangential_projected} route.

Let the shell electric field be split into a known driving part and the
response associated with the unknown toroidal update:
\[
\mathbf E = \mathbf E^{\rm drv} + \mathbf E^{\rm ind}[\partial_t\psi].
\]

There are then two exact shell closures one may write.

\subsection*{Exact radial shell closure}

From the radial component of
\[
\mu_0 \partial_t \mathbf j = -\nabla\times\nabla\times \mathbf E,
\]
one gets the exact shell equation
\begin{equation}
\mu_0 B_{0r}\eta
=
-\frac{1}{R_I^2}\Bigl[Q_I - \Delta_\Omega E_{r,I}\Bigr],
\label{eq:radial_shell_balance}
\end{equation}
with
\[
Q_I := \partial_r\!\bigl(r\,\mathrm{div}_\Omega \mathbf E_S\bigr)\Big|_{R_I^+}.
\]
This is the closure used in the earlier note.

\subsection*{Exact tangentially projected shell closure}

From the tangential part of the same Maxwell equation,
\[
\mu_0 \eta \mathbf B_{0S} = \bigl[-\nabla\times\nabla\times \mathbf E\bigr]_S.
\]
Projecting onto $\mathbf B_{0S}$ gives
\begin{equation}
\mu_0 |\mathbf B_{0S}|^2 \eta
=
\mathbf B_{0S}\cdot \bigl[-\nabla\times\nabla\times \mathbf E\bigr]_S.
\label{eq:tangential_projected_driver}
\end{equation}
If the induced part is rewritten through \eqref{eq:tangential_projected_balance_no_driver},
the closed driven equation becomes
\begin{equation}
\mu_0 |\mathbf B_{0S}|^2 \eta
=
f_t^{\rm drv}
+
\mathbf B_{0S}\cdot \nabla_S\!\bigl(\partial_t\psi + R_I \partial_r\partial_t\psi\bigr),
\label{eq:tangential_projected_balance}
\end{equation}
where
\begin{equation}
f_t^{\rm drv}
:=
\mathbf B_{0S}\cdot \bigl[-\nabla\times\nabla\times \mathbf E^{\rm drv}\bigr]_S.
\label{eq:f_t_drv}
\end{equation}

Equation \eqref{eq:tangential_projected_balance} is the exact continuous
projected shell equation that is structurally closest to the present
\texttt{dt\_alpha} formulation. But it is still only a projected scalar closure,
not the full tangential vector balance.

\section{The first-principles projected operator}

Using \eqref{eq:dtpsi_from_eta}, \eqref{eq:dr_dtjr_unit_sphere},
\eqref{eq:dr_dtpsi}, and writing
\[
A_\Omega := \mathbf B_{0S}\cdot \nabla_\Omega = R_I\,\mathbf B_{0S}\cdot \nabla_S,
\]
equation \eqref{eq:tangential_projected_balance} becomes
\begin{equation}
\mu_0 |\mathbf B_{0S}|^2 \eta
=
f_t^{\rm drv}
+
A_\Omega\!\left[
\frac{1}{R_I}\,\partial_t\psi + \partial_r\partial_t\psi
\right].
\label{eq:code_matching_continuous}
\end{equation}

Introduce the exact continuous maps
\[
T_{\alpha\to j_r}\eta = B_{0r}\eta,
\qquad
T_{j_r\to\psi} = -\mu_0 R_I \Delta_\Omega^{-1},
\qquad
T_{\alpha\to\psi} = T_{j_r\to\psi} T_{\alpha\to j_r}.
\]
Also define
\[
D_r^{j_r\leftarrow\alpha}\eta
:=
(\partial_r B_{0r})\eta - \frac{1}{R_I}A_\Omega \eta,
\]
and
\[
D_r^{\psi\leftarrow\alpha}
:=
\frac{1}{R_I}T_{\alpha\to\psi}
+
T_{j_r\to\psi} D_r^{j_r\leftarrow\alpha}.
\]
Then the exact tangential shell operator can be written
\begin{equation}
L_\alpha \eta
:=
\mu_0 |\mathbf B_{0S}|^2 \eta
-\,
A_\Omega\!\left[
\frac{1}{R_I} T_{\alpha\to\psi}
+
D_r^{\psi\leftarrow\alpha}
\right]\eta.
\label{eq:L_alpha_continuous}
\end{equation}

The exact projected shell closure is therefore
\begin{equation}
L_\alpha \eta = f_t^{\rm drv}.
\label{eq:code_matching_closure}
\end{equation}

\paragraph{The dropped perpendicular component.}
Let
\[
\mathbf B_{0S}^{\perp} := (-B_{0\phi},\,B_{0\theta})
\]
denote the in-shell vector perpendicular to $\mathbf B_{0S}$. Then the part of
the full tangential balance discarded by the scalar projection is
\begin{equation}
R_\perp(\eta)
:=
\mathbf B_{0S}^{\perp}\cdot
\left[
\mu_0 \eta \mathbf B_{0S}
-
\nabla_S\!\bigl(\partial_t\psi + R_I \partial_r\partial_t\psi\bigr)
\right].
\label{eq:R_perp}
\end{equation}
The projected closure keeps only
\[
\mathbf B_{0S}\cdot[\cdots],
\]
while $R_\perp$ is unconstrained. This is the precise first-principles quantity
that should be monitored if one wants to assess how much tangential Maxwell/FAC
content is being discarded by the projected scalar route.

\paragraph{Equivalent feedback form.}
Since
\[
D_r^{\psi\leftarrow\alpha}
=
\frac{1}{R_I}T_{\alpha\to\psi}
+
T_{j_r\to\psi} D_r^{j_r\leftarrow\alpha},
\]
the feedback block may also be written
\[
A_\Omega\!\left[
\frac{2}{R_I}T_{\alpha\to\psi}
+
T_{j_r\to\psi} D_r^{j_r\leftarrow\alpha}
\right]
\]
with an overall minus sign in \eqref{eq:L_alpha_continuous}. Equivalently,
\[
L_\alpha
=
\mu_0 |\mathbf B_{0S}|^2 I
-\,
A_\Omega\!\left[
\frac{2}{R_I}T_{\alpha\to\psi}
+
T_{j_r\to\psi} D_r^{j_r\leftarrow\alpha}
\right].
\]

\section{Comparison to the current code operator}

\paragraph{Scope of this sign audit.}
This section and the next two concern only the
\path{tangential_projected} branch. They do \emph{not} diagnose the newer
\path{radial_shell} runtime directly. Their role is narrower: to show which
parts of the older projected tangential operator are exact shell identities,
and which parts still require an independent sign or closure justification.

The current implementation in \texttt{toroidal.py} assembles the operator
\[
L_\alpha^{\rm code}
=
\mu_0 |\mathbf B_{0S}|^2 I
+
A_\Omega\!\left[
\frac{2}{R_I}T_{\alpha\to\psi}
+
T_{j_r\to\psi} D_r^{j_r\leftarrow\alpha}
\right],
\]
that is, with the \emph{opposite sign} on the feedback block from the
first-principles projected shell operator \eqref{eq:L_alpha_continuous}, under
the conventions and operator definitions used in this note.

So the present note does \emph{not} prove that the current assembled sign is
correct. What it proves is narrower:
\begin{itemize}
\item the chain
$\eta \to \partial_t j_r \to \partial_t\psi \to \partial_r\partial_t\psi$ is
first-principles;
\item the appearance of $|\mathbf B_{0S}|^2$ belongs naturally to a projected
tangential shell equation;
\item but the sign of the assembled feedback block in the current code remains
an unresolved issue and requires its own audit.
\end{itemize}

\section{What is exact in the current implementation}

The following \emph{building blocks} of the current toroidal full-induction path
are derivable from the assumptions above:
\begin{enumerate}
\item $T_{\alpha\to j_r}$ from $j_r = B_{0r}\alpha$;
\item $T_{j_r\to\psi}$ from
$j_r = -(\mu_0 R_I)^{-1}\Delta_\Omega \psi$;
\item $T_{\alpha\to\psi} = T_{j_r\to\psi} T_{\alpha\to j_r}$;
\item $D_r^{j_r\leftarrow\alpha}$ from field-line invariance and
$\nabla\cdot \mathbf B_0 = 0$;
\item $D_r^{\psi\leftarrow\alpha}$ from differentiating the shell relation for
$\psi$;
\item $A_\Omega = \mathbf B_{0S}\cdot \nabla_\Omega$;
\item the use of $|\mathbf B_{0S}|^2$ in the tangentially projected shell
balance.
\end{enumerate}

So the present \texttt{dt\_alpha} path contains several exact structural
ingredients. But the \emph{assembled sign} of the feedback block is not shown
correct by this note; at present it appears opposite to the sign obtained from
the projected tangential shell balance.

\section{What still requires explicit justification}

The following parts of the current implementation are not pure consequences of
the basic assumptions and should be documented as additional closures or
numerical choices:

\begin{enumerate}
\item \textbf{The explicit forcing operator
$\mathbf E^{\rm drv} \mapsto f_t^{\rm drv}$.}

The code uses an Er-free shell operator built from the tangential field:
\[
c = \mathrm{curl}_\Omega(\mathbf E_S),
\qquad
f_t^{\rm drv}
\sim
\frac{1}{R_I^2}
\left(
B_{0\theta}\partial_\phi^\ast c - B_{0\phi}\partial_\theta c
\right).
\]
But the sign is not fixed by the present derivation. Under the natural
shell-local extension
\[
E_r = 0,
\qquad
\partial_r \mathbf E_S = 0
\quad\text{at } r=R_I^+,
\]
the direct tangential shell trace gives
\[
\mathbf B_{0S}\cdot[-\nabla\times\nabla\times \mathbf E]_S
=
-\,\frac{1}{R_I^2}
\left(
B_{0\theta}\partial_\phi^\ast c - B_{0\phi}\partial_\theta c
\right),
\]
that is, the \emph{opposite sign} from the currently assembled code operator.

So the present Er-free forcing block should not yet be described as
first-principles. A complete derivation must show exactly how the eliminated
radial-electric terms change this shell-local sign, or else the implementation
sign should be treated as suspect.

\item \textbf{The projection from the full tangential vector equation to a
scalar closure.}

The projected equation keeps only the component parallel to $\mathbf B_{0S}$.
That is an exact reduction, but it is still a weaker equation than the full
tangential Maxwell/FAC balance.

\item \textbf{The inverse map $j_r \mapsto \alpha$.}

The code uses a regularized weak-form minimum-norm branch selector near small
$|B_{0r}|$. This is a numerical regularization, not an exact identity.

\item \textbf{Hard and soft least-squares constraints.}

The weighted least-squares solve, soft penalties, and Tikhonov regularization
are numerical choices. They are reasonable, but they are not implied by the
continuous shell equations.

\item \textbf{Hemisphere coupling.}

The current \texttt{connect\_hemispheres=True} implementation uses compressed
low-latitude apex mismatch constraints. That is not the same as an exact
conjugate-footpoint equality constraint
\[
\eta(\Omega)=\eta(\mathcal C(\Omega)).
\]
It should therefore be presented as a model surrogate, not as the exact closed
field-line topology condition.

\item \textbf{Auxiliary closure basis choices.}

Using an auxiliary SH closure basis for CS-dominant full induction is an
implementation architecture choice. It may be good numerically, but it is not a
first-principles statement.

\item \textbf{The full upper-region induced response.}

This note derives only the local toroidal/FAC shell backbone. A complete
full-induction closure note still needs the nonlocal upper-region response
operator that maps the shell update to the full induced electric and magnetic
response above the shell.
\end{enumerate}

\section{Radial-shell closure and the missing scalar}

The exact radial shell closure is
\begin{equation}
\mu_0 B_{0r}\eta
=
-\frac{1}{R_I^2}\Bigl[Q_I - \Delta_\Omega E_{r,I}\Bigr].
\label{eq:radial_shell_balance_repeat}
\end{equation}
The unresolved upper-side quantity can be written in several exactly
equivalent ways. From the upper-side toroidal relation,
\begin{equation}
\mu_0\,\partial_t j_r^+
=
-\,\frac{1}{R_I}\Delta_\Omega \partial_t\psi^+,
\label{eq:dtjr_dtpsi_upper}
\end{equation}
while the radial Maxwell balance gives
\begin{equation}
\mu_0\,\partial_t j_r^+
=
-\,\frac{1}{R_I^2}\Bigl[Q_I - \Delta_\Omega E_{r,I}\Bigr].
\label{eq:dtjr_electric_upper}
\end{equation}
Therefore
\begin{equation}
-\,\frac{1}{R_I^2}\Bigl[Q_I - \Delta_\Omega E_{r,I}\Bigr]
\quad\Longleftrightarrow\quad
\mu_0\,\partial_t j_r^+
\quad\Longleftrightarrow\quad
-\,R_I \Delta_S(\partial_t\psi^+).
\label{eq:missing_scalar_equivalence}
\end{equation}

\paragraph{Direct derivation from the electric decomposition.}
Using
\[
\mathbf E
=
E_r\,\hat r + \nabla_S U - \hat r\times \nabla_S V,
\]
the first curl splits as
\[
\nabla\times \mathbf E
=
\omega_E\,\hat r
+
\hat r\times \nabla_S(\partial_r U - E_r)
+
\nabla_S(\partial_r V),
\]
where
\[
\omega_E = \hat r\cdot(\nabla\times \mathbf E).
\]
The radial component of the second curl depends only on the tangential part of
\(\nabla\times \mathbf E\), and the gradient term \(\nabla_S(\partial_r V)\)
drops out of the radial curl. Therefore
\begin{equation}
\bigl[-\nabla\times\nabla\times \mathbf E\bigr]_r
=
-\,\Delta_S(\partial_r U - E_r)
=
-\,\frac{1}{R_I^2}\Bigl[\Delta_\Omega(\partial_r U) - \Delta_\Omega E_r\Bigr].
\label{eq:radial_shell_from_first_traces}
\end{equation}
Since
\[
Q_I = \Delta_\Omega(\partial_r U)\big|_{R_I^+},
\]
equation \eqref{eq:radial_shell_from_first_traces} is exactly the radial shell
closure \eqref{eq:radial_shell_balance_repeat}. So the exact radial-shell
scalar depends only on the \emph{first} shell traces
\[
\partial_r U\big|_{R_I^+},
\qquad
E_{r,I}.
\]
This is the clean contrast with the omitted projected-tangential forcing term,
which requires second radial trace information.

\paragraph{Meaning.}
The radial-shell closure is not missing ``the whole off-shell electric field''.
It is missing one scalar upper-side update quantity. One may obtain it through:
\begin{itemize}
\item an electric-trace route
\[
\partial_t\psi \mapsto \bigl(E_{r,I},\,Q_I\bigr),
\]
or
\item a toroidal-update route
\[
\partial_t\psi \mapsto \partial_t\psi^+
\quad\text{or}\quad
\partial_t j_r^+.
\]
\end{itemize}
Either closes the same exact scalar equation through
\eqref{eq:missing_scalar_equivalence}.

\paragraph{One-curl versus two-curl viewpoints.}
There are two equally exact ways to reach the same scalar closure from the
electric field just above the shell.
\begin{enumerate}
\item \textbf{One-curl / connector-scalar route.} The toroidal part of
Faraday's law gives
\[
R_I\,\partial_t\psi^+
=
\Pi_0(\partial_r U - E_r).
\]
This is the cleanest route if the primary missing object is taken to be the
upper-side toroidal update \(\partial_t\psi^+\).

\item \textbf{Two-curl / radial-current route.} Taking the radial component of
\(-\nabla\times\nabla\times \mathbf E\) gives
\[
\mu_0\,\partial_t j_r^+
=
-\,\Delta_S(\partial_r U - E_r)
=
-\,R_I\Delta_S(\partial_t\psi^+).
\]
This is the cleanest route if the primary missing object is taken to be the
upper-side radial-current update \(\partial_t j_r^+\).
\end{enumerate}
So the two-curl route is not adding new physics; it is the radial-current form
of the same scalar connector relation. The one-curl route is usually more
intuitive for \(\partial_t\psi^+\), while the two-curl route is usually more
intuitive for the radial-shell balance itself.

\paragraph{Old ``surface'' versus ``radial-structure'' language.}
The older note sometimes described the radial-shell forcing in terms of a
``surface part'' and a ``radial-structure part.'' In the present formulation
that language should be used carefully:
\begin{itemize}
\item for the \emph{exact radial-shell connector}, the cleaner object is the
single scalar
\[
\partial_r U - E_r
\quad\Longleftrightarrow\quad
R_I\,\partial_t\psi^+.
\]
So on the exact radial-shell side it is better to speak directly about the
connector scalar or the first traces \((\partial_r U, E_r)\), rather than a
surface-versus-radial split;

\item for the \emph{reduced projected-tangential forcing} route, the old
language still applies. There the live operator keeps the shell-only
``surface'' contribution coming from the radial curl scalar
\[
\omega_E = \hat r\cdot(\nabla\times \mathbf E),
\]
while the omitted complementary piece depends on explicit off-shell radial
trace information.
\end{itemize}
So the newer one-curl framework does not invalidate the old distinction; it
simply shows that the distinction belongs mainly to the reduced forcing-side
decomposition, not to the exact radial-shell connector itself.

\section{A first-principles radial-shell response operator}

To turn \eqref{eq:radial_shell_balance_repeat} into a closed equation in
$\eta$, one needs the shell scalar response produced by a trial toroidal shell
update
\[
\partial_t\psi = -\,\mu_0 R_I\,\Delta_\Omega^{-1}\Pi_0(B_{0r}\eta).
\]

The strongest first-principles upper-region model compatible with the present
code assumptions is:
\begin{enumerate}
\item quasi-static Maxwell equations in the gap;
\item strong field-aligned current in the gap, written with an upper-region FAC
amplitude \(\beta\),
\[
\partial_t \mathbf j^+ = \beta\,\mathbf B_0;
\]
\item current continuity in the gap, hence
\[
\mathbf B_0\cdot\nabla\beta=0;
\]
\item the shell trial variable \(\eta\) supplies the lower-boundary/interface
data for \(\beta\); on closed lines this means \(\beta\) agrees with the
field-line continuation of the shell trial value, while open lines use the
chosen outer semantics;
\item an explicit open-line / outer-boundary condition at $R_M$ or infinity;
\item the ideal electrodynamic closure in the gap
\[
\mathbf E\cdot\mathbf B_0 = 0,
\]
\item a charge-free bulk-gap closure for the perturbation electric field
\[
\nabla\cdot \mathbf E^+ = 0,
\]
which is adopted for the region above the ionosphere and should be read as part
of the bulk-gap model, not as part of the minimal exact shell backbone.
\end{enumerate}

\paragraph{Meaning of quasi-static Maxwell here.}
In this note, ``quasi-static Maxwell'' means that the gap obeys
\[
\nabla\times \mathbf B = \mu_0 \mathbf j,
\qquad
\nabla\times \mathbf E = -\,\partial_t\mathbf B,
\qquad
\nabla\cdot \mathbf B = 0,
\]
with displacement current neglected. So induction is retained through Faraday's
law, but electromagnetic wave propagation and retardation are not part of the
model.

With those assumptions, a trial $\eta$ defines a trial $\partial_t\psi$ on the
shell and therefore the lower-boundary/interface data for the gap FAC
amplitude \(\beta\) and the associated toroidal magnetic update in the gap.
The exact nonlocal radial-shell response operator is then
\begin{equation}
\mathcal R_r[\partial_t\psi]
:=
-\frac{1}{R_I^2}
\left[
Q_I^{\rm ind}[\partial_t\psi]
- \Delta_\Omega E_{r,I}^{\rm ind}[\partial_t\psi]
\right],
\label{eq:radial_response_operator}
\end{equation}
or, equivalently, the upper-side toroidal update operator
\[
\partial_t\psi \mapsto \partial_t\psi^+
\quad\text{or}\quad
\partial_t\psi \mapsto \partial_t j_r^+.
\]

If
\[
f_r^{\rm drv}
:=
-\frac{1}{R_I^2}
\left[
Q_I^{\rm drv} - \Delta_\Omega E_{r,I}^{\rm drv}
\right],
\]
then the exact closed radial-shell problem is
\begin{equation}
\mu_0 B_{0r}\eta
=
f_r^{\rm drv}
+
\mathcal R_r[\mathcal A\eta],
\qquad
\mathcal A\eta
:=
-\,\mu_0 R_I\,\Delta_\Omega^{-1}\Pi_0(B_{0r}\eta).
\label{eq:closed_radial_shell_response_problem}
\end{equation}

\paragraph{This remains the canonical scalar closure.}
Equation \eqref{eq:closed_radial_shell_response_problem} is the cleanest
scalar first-principles statement of the toroidal induction closure. No
projection of a tangential vector equation is required.

\paragraph{Exact gap-BVP viewpoint.}
The strongest independent upper-region closure suggested by the present model
is not a new local shell law. It is a global quasi-static ideal
boundary-value problem in the gap, whose shell condensation yields the missing
nonlocal operator.

\paragraph{Two related gap problems.}
There are in fact two related upper-region problems in this note.

\emph{Gap problem A: induced-response operator.}
A trial shell toroidal update \(\partial_t\psi\) (equivalently a shell trial
\(\eta\)) is prescribed, and the gap returns the induced radial-shell response
operator \(\mathcal R_r\), or equivalently the induced \(\partial_t\psi^+\) or
\(\partial_t j_r^+\).

\emph{Gap problem B: forcing-side DtN operator.}
The shell tangential electric trace \(\mathbf E_{S,\mathrm{shell}}\) is
prescribed, and the gap returns the condensed source operator
\(\Gamma_{\rm known}\), or equivalently the trace map \(\mathcal D_q\), the
connector scalar \(\chi\), or the source \(s^+=\partial_t j_r^+\).

The same upper-region model underlies both, but their prescribed shell data
and returned shell objects are different.

\paragraph{Which one is the real unresolved closure?}
At the level of the reduced-shell theory, the primary closure object is still
\(\Pi=\mathcal P_I[\chi]\). But at the level of \emph{constructive runtime
assembly}, Gap problem B is the main unresolved forcing-side interface:
once an upper source distribution, shell correction, or returned toroidal
trace is specified above the ionosphere, the subsequent magnetic continuation
through the gap is already comparatively routine in the existing machinery.
So the main remaining constructive connector problem is
\[
\text{ionosphere shell state}
\;\to\;
q=\partial_rU-E_r
\;\to\;
\chi=\partial_t\psi^+,
\]
or equivalently the returned upper source \(s^+=\partial_t j_r^+\). Gap
problem A remains useful as an induced-branch decomposition and diagnostic,
but it is not the primary unresolved forcing-side object.

Let the upper region be
\[
\mathcal G=\{R_I<r<R_M\}
\]
or \(R_I<r<\infty\) in the open case. Introduce an upper-region FAC amplitude
\(\beta\) such that
\[
\partial_t\mathbf j^+ = \beta\,\mathbf B_0,
\qquad
\mathbf B_0\cdot\nabla\beta = 0.
\]
Then the minimal adopted upper-region system is
\begin{align}
\nabla\times \partial_t\mathbf B^+ &= \mu_0\,\beta\,\mathbf B_0,
\label{eq:gap_bvp_ampere}
\\
\nabla\cdot \partial_t\mathbf B^+ &= 0,
\label{eq:gap_bvp_divb}
\\
\nabla\times \mathbf E^+ &= -\,\partial_t\mathbf B^+,
\label{eq:gap_bvp_faraday}
\\
\mathbf E^+\cdot \mathbf B_0 &= 0.
\label{eq:gap_bvp_ideality}
\\
\nabla\cdot \mathbf E^+ &= 0.
\label{eq:gap_bvp_dive}
\end{align}
This is the quasi-static, ideal, charge-free, strong-FAC gap model in its most
compact form. Here \(\beta\) is not a freely chosen source term; it is part of
the upper-region solution and must be tied to the shell interface/current-balance
conditions together with the chosen open/closed-line semantics. In the
induced-response problem, \(\beta\) is the gap continuation of the prescribed
shell trial \(\eta\). In the forcing-side closure problem, \(\beta\) is solved
as part of the gap state consistent with the prescribed shell tangential
electric trace and the adopted boundary/interface conditions. The condition
\(\nabla\cdot \mathbf E^+=0\) is an adopted bulk-gap closure assumption for the
region above the ionosphere; it is useful because it makes the curl-free
remainder of the electric field a legitimate harmonic-continuation target, but
it should not be confused with the exact shell identities derived earlier in
the note.

The upper-region problem is not closed until the boundary and interface data
are fixed. The minimal required choices are:
\begin{enumerate}
\item the shell interface data, chosen according to which gap problem is being
posed: for the forcing-side closure problem one prescribes the upper
tangential electric trace
\[
\mathbf E_S^+\big|_{S_I}=\mathbf E_{S,\mathrm{shell}},
\]
and solves for the outgoing upper-side connector response
\[
\chi=\partial_t\psi^+\big|_{S_I}
\quad\text{or equivalently}\quad
\partial_t j_r^+\big|_{S_I};
\]
for the induced-response problem, the natural shell input is the trial update
\(\eta\) or equivalently \(\partial_t\psi\); if the gap is coupled back to the
ionosphere as part of a full shell-gap system, these interface data must be
paired with the adopted sheet-jump/current-balance law at the ionosphere;
\item the outer boundary condition. For the model variants actually targeted in
this note, open lines are closed by \emph{one of two admissible outer
semantics}: decay to infinity or an \(R_M\)-shielded response. More general
outer conditions may be posed in principle, but they are not the present
closure target. For the finite-\(R_M\) variant, the cleanest explicit outer
magnetic boundary pair is
\[
\psi^+(R_M),
\qquad
M(R_M):=\bigl(-R_M^2\Delta_S\bigr)^\dagger B_r(R_M),
\]
that is, the toroidal magnetic scalar together with the poloidal magnetic
scalar reconstructed from the radial field at \(R_M\);
\item the closed-line/open-line semantics for \(\beta\) on the chosen field
topology;
\item the mean-zero/solvability convention used to invert
\(\Delta_\Omega\).
\end{enumerate}

\paragraph{Interface unknowns in operator form.}
For the forcing-side gap problem alone, the clean mixed formulation is:
\[
\mathbf E_S^+\big|_{S_I}=\mathbf E_{S,\mathrm{shell}}
\quad\text{prescribed},
\]
solve in the gap for
\[
(\mathbf E^+,\partial_t\mathbf B^+,\beta),
\]
and then define the outgoing upper-side connector scalar and source through
\[
\chi=\frac{1}{R_I}\Pi_0(\partial_r U-E_r),
\qquad
\mu_0\,\partial_t j_r^+ = -\,R_I\Delta_S\chi.
\]
This is the precise shell/gap interface viewpoint behind
\(\Gamma_{\rm known}\) and \(\mathcal D_q\): the shell prescribes the
tangential electric trace, while the gap solve returns the connector scalar or
equivalently the outgoing upper radial-current source.

\paragraph{Why the connector scalar is the right unknown.}
The residual potential freedom compatible with ideality does not change the
connector scalar. Indeed, if
\[
\mathbf E^+ \mapsto \mathbf E^+ + \nabla\phi,
\qquad
\mathbf B_0\cdot\nabla\phi = 0,
\]
then \(U\mapsto U+\phi\) and \(E_r\mapsto E_r+\partial_r\phi\), so the
combination
\begin{equation}
q := \partial_r U - E_r
\label{eq:q_gap_bvp}
\end{equation}
is invariant. Therefore the shell connector quantity
\[
\chi := \frac{1}{R_I}\Pi_0 q
\]
is the natural shell variable even if \(\mathbf E^+\) itself retains a
gauge-like freedom before boundary conditions are imposed.
Because \(U\) here already includes the repo curl-free Helmholtz sign,
\eqref{eq:q_gap_bvp} is exactly the legacy-sign-compatible connector scalar
used by the current implementation. If one writes the shell field directly in
terms of the raw Helmholtz coefficient \(\Phi\), then
\[
q = s_{\rm cf}\,\partial_r \Phi - E_r.
\]

\paragraph{Connector extraction from the gap solution.}
With
\[
\mathbf E^+
=
E_r\,\hat r + \nabla_S U - \hat r\times \nabla_S V,
\]
the first curl takes the form
\[
\nabla\times \mathbf E^+
=
\omega_E\,\hat r
+
\hat r\times \nabla_S q
+
\nabla_S p,
\qquad
p:=\partial_r V.
\]
Hence the upper-side tangential magnetic trace contains the connector through
\[
(\partial_t\mathbf B^+)_S
=
-\,\hat r\times \nabla_S q
- \nabla_S p.
\]
So the gap BVP does not have to guess the missing scalar; it produces
\(\chi\) directly through the \(\hat r\times \nabla_S q\) part of the
tangential magnetic trace. Equivalently,
\[
\mu_0\,\partial_t j_r^+
=
-\,R_I\Delta_S \chi.
\]

\paragraph{What the bulk charge-free closure fixes locally.}
The adopted bulk condition \(\nabla\cdot \mathbf E^+=0\) gives a genuine
continuation law for the curl-free tangential electric part. With
\[
\mathbf E^+
=
E_r\,\hat r + \nabla_S U - \hat r\times \nabla_S V,
\]
one has
\begin{equation}
\nabla\cdot \mathbf E^+
=
\frac{1}{r^2}\partial_r(r^2 E_r) + \Delta_S U
= 0.
\label{eq:gap_divfree_E}
\end{equation}
So at the shell,
\begin{equation}
\partial_r E_r\big|_{R_I^+}
=
-\,\frac{2}{R_I}E_{r,I} - \Delta_S U.
\label{eq:gap_divfree_E_shell}
\end{equation}
This removes arbitrary bulk charge freedom, but it still does not determine
the connector scalar \(q=\partial_r U-E_r\) by itself.

\paragraph{Differentiated ideality and the remaining local freedom.}
With
\[
\mathbf E^+\cdot\mathbf B_0 = 0,
\]
the shell decomposition gives
\[
B_{0r}E_r + \mathbf B_{0S}\cdot\nabla_S U - \mathbf B_{0S}^{\perp}\cdot\nabla_S V = 0.
\]
Differentiating radially and writing
\[
q:=\partial_r U - E_r,
\qquad
p:=\partial_r V,
\]
one obtains, after using \eqref{eq:gap_divfree_E_shell},
\begin{equation}
\mathbf B_{0S}\cdot\nabla_S q
- \mathbf B_{0S}^{\perp}\cdot\nabla_S p
=
F_q[U,V,E_r;\mathbf B_0],
\label{eq:q_p_surface_relation}
\end{equation}
where
\begin{align}
F_q[U,V,E_r;\mathbf B_0]
:={}&
-\,\mathbf B_{0S}\cdot\nabla_S E_r
+ B_{0r}\!\left(\frac{2}{R_I}E_r + \Delta_S U\right)
- (\partial_r B_{0r})E_r
\notag\\
&-\,(\partial_r \mathbf B_{0S})\cdot\nabla_S U
+ (\partial_r \mathbf B_{0S}^{\perp})\cdot\nabla_S V.
\label{eq:q_p_surface_rhs}
\end{align}
For fixed shell \(U\), \(V\), and \(E_r\) (with \(E_r\) itself determined
away from \(B_{0r}=0\) by shell ideality), the right-hand side is known from shell
data and the background-field radial variation. So the bulk equations no
longer leave arbitrary radial structure: locally they reduce the unresolved
first traces to the pair \((q,p)\) constrained by one scalar surface relation.
The remaining freedom is therefore not another missing local shell identity,
but one genuinely global boundary/topology degree selected by the chosen gap
boundary-value problem.

\paragraph{Condensed shell operator.}
If the adopted gap BVP is well posed, its shell condensation defines a
Dirichlet-to-Neumann-type map of the form
\[
\chi = \mathcal T_{\rm gap}\big[\mathbf E_{S,\mathrm{shell}}\big],
\]
or equivalently
\[
-\,\frac{R_I}{\mu_0}\Delta_S\chi
=
\Gamma_{\mathrm{known}}\,\mathbf E_{S,\mathrm{shell}}.
\]
In this language, \(\Gamma_{\rm known}\) is no longer a surrogate forcing
block. It is the shell operator that the adopted gap BVP would define after
condensation.

\paragraph{Scope of this claim.}
This is the strongest formulation the present note can defend. It does
\emph{not} yet prove well-posedness of the gap BVP or construct a closed-form
matrix for \(\Gamma_{\rm known}\). What it does establish is the correct
mathematical target: the independent forcing-side closure, if obtained, should
come from solving and condensing the global upper-region boundary-value
problem, not from guessing an additional local shell scalar law. So the
closure is now complete in structural/operator form, but not yet in a fully
constructive assembled form for one chosen discretized gap model. In that
sense, the local shell algebra by itself remains underdetermined, while the
shell-plus-chosen-gap-BVP problem is the actual solvable object.

\paragraph{Abstract constrained co-energy closure.}
If one wants a shell-only operator for the missing forcing-side closure, the
most natural abstract candidate is a constrained magnetic co-energy
minimization in the gap. Carry the prescribed shell data explicitly by fixing
\[
\mathbf E_{S,\rm known}
\quad\text{and}\quad
\chi:=\partial_t\psi^+\big|_{S_I},
\]
and let \(\mathcal A(\mathbf E_{S,\rm known},\chi)\) denote the admissible gap
fields satisfying the same quasi-static, ideal, strong-FAC, closed/open-line,
and outer-boundary assumptions used above. Then the condensed shell functional
is
\begin{equation}
\mathcal J_{\rm eff}[\chi;\mathbf E_{S,\rm known}]
:=
\inf_{X\in\mathcal A(\mathbf E_{S,\rm known},\chi)}
\frac{1}{2\mu_0}\int_{\mathcal G} |\partial_t\mathbf B[X]|^2\,dV.
\label{eq:abstract_gap_coenergy_functional}
\end{equation}
After eliminating the gap variables, one aims to obtain the shell form
\begin{equation}
\mathcal J_{\rm eff}[\chi;\mathbf E_{S,\rm known}]
=
\frac12 \langle \chi,\Lambda_{\rm gap}\chi\rangle
- \langle \Gamma_{\rm known}\mathbf E_{S,\rm known}, \chi\rangle,
\label{eq:abstract_shell_coenergy}
\end{equation}
with Euler--Lagrange equation
\begin{equation}
\Lambda_{\rm gap}\chi = \Gamma_{\rm known}\mathbf E_{S,\rm known}.
\label{eq:abstract_shell_coenergy_EL}
\end{equation}
The special case with a generic shell/source term \(g_{\rm surf}\) is just
\eqref{eq:abstract_shell_coenergy} with
\[
g_{\rm surf}=\Gamma_{\rm known}\mathbf E_{S,\rm known},
\]
so the older ``surface term'' versus ``radial-structure'' split becomes:
\[
\text{surface/source term } g_{\rm surf},
\qquad
\text{radial structure } \Lambda_{\rm gap}.
\]

\paragraph{What is exact already, and what is still nonlocal.}
In the variable choice
\[
\chi=\partial_t\psi^+,
\]
the final shell inversion is already exact:
\[
\mu_0\,\partial_t j_r^+ = -\,R_I\Delta_S\chi,
\qquad
\Lambda_{\rm shell}^{\rm exact}:= -\,\frac{R_I}{\mu_0}\Delta_S.
\label{eq:exact_shell_radial_structure_operator}
\]
So the genuinely nonlocal object is not the final radial structure itself, but
the source builder \(\Gamma_{\rm known}\) determined by the condensed gap
model.

\paragraph{Reduction to the q-trace map.}
With ideality adopted, the known shell tangential data already determine
\[
E_{r,\rm known}
=
-\,\frac{\mathbf B_{0S}\cdot \mathbf E_{S,\rm known}}{B_{0r}},
\qquad
(B_{0r}\neq 0),
\]
so the only missing first trace is
\[
q_{\rm known}:=\partial_r U_{\rm known}-E_{r,\rm known}.
\]
Near small \(B_{0r}\), this relation should be understood through the
shell-source/connector formulation rather than as a pointwise division rule.
The one-curl identity then gives
\begin{equation}
R_I\,\chi_{\rm known} = \Pi_0 q_{\rm known},
\qquad
s_{\rm known}^+
=
-\,\frac{1}{\mu_0 R_I^2}\,\Delta_\Omega q_{\rm known}.
\label{eq:gamma_known_from_q}
\end{equation}
Hence
\begin{equation}
\Gamma_{\rm known}
=
-\,\frac{1}{\mu_0 R_I^2}\,\Delta_\Omega\,\mathcal D_q,
\label{eq:gamma_known_via_dtn_q}
\end{equation}
where
\[
\mathcal D_q:\mathbf E_{S,\rm known}\mapsto q_{\rm known}
\]
is the smallest exact forcing-side interface. Conversely, whenever a forcing
model already provides \(\Gamma_{\rm known}\), the corresponding \(q\)-trace
follows by exact shell inversion:
\begin{equation}
\mathcal D_q
=
R_I\,T_{j_r\to\psi}\,\Gamma_{\rm known}
=
-\,\mu_0 R_I^2\,\Delta_\Omega^{-1}\Gamma_{\rm known}.
\label{eq:Dq_from_source_operator}
\end{equation}
This is why the forcing-side \emph{construction} is most conveniently organized
around \(q\), and why the concrete \(q\)-trace realizations are deferred to
the runtime appendix. But the shell-facing unknown remains
\[
\chi=\frac{1}{R_I}\Pi_0 q,
\]
so \(q\) should be read as the preimage used to build the returned shell
variable, not as a competing primary shell unknown.

\paragraph{Relation to the reduced-shell magnetic operator.}
The operators \(\mathcal D_q\) and \(\Gamma_{\rm known}\) belong to the
\emph{forcing-side/electric-side} version of the closure problem: they start
from prescribed shell forcing data and return the connector scalar or the
equivalent shell source. The reduced-shell formulation of the companion
document is different in emphasis. There the primitive unknown is already the
shell toroidal update, which that document denotes by \(q\). In the notation
of the present note, that reduced-shell unknown is
\[
\chi=\partial_t\psi^+ = \frac{1}{R_I}\Pi_0 q,
\]
the harmonic baseline fixed by the shell radial magnetic update is subtracted,
and the remaining tangential poloidal magnetic correction is written
\[
\Pi = \mathcal P_I[\chi].
\]
So \(\mathcal D_q\), \(\Gamma_{\rm known}\), and \(\mathcal P_I\) should be
read as three faces of the same upper response:
\begin{itemize}
\item \(\mathcal D_q\) is the shell-electric to connector map;
\item \(\Gamma_{\rm known}\) is the corresponding shell source builder;
\item \(\mathcal P_I\) is the shell-magnetic correction operator once
\(\chi\) is already the reduced-shell unknown.
\end{itemize}
This is the main conceptual bridge between the present forcing-side
construction and the compressed shell closure
\(K_I \chi = -R_I\widetilde G\).

\paragraph{Minimal direct closure.}
The least messy explicit forcing-side model is therefore to introduce one
side-trace operator
\[
\Lambda_U: U|_{S_I}\mapsto \partial_r U|_{R_I^+},
\]
and then set
\[
q_{\rm known} = \Lambda_U U - E_{r,\rm known},
\qquad
\chi_{\rm known}=\frac{1}{R_I}\Pi_0 q_{\rm known}.
\]
This isolates the only genuinely nonlocal ingredient into a single shell-side
matrix. Open harmonic and \(R_M\)-shielded harmonic Dirichlet-to-Neumann maps
for \(U\) are the simplest explicit realizations of \(\Lambda_U\).

\paragraph{Particular plus homogeneous refinement.}
One need not assume that the whole forcing-side connector is harmonic. A more
refined split is
\[
\mathbf E = \mathbf E_{\rm part} + \mathbf E_{\rm hom},
\qquad
\nabla\times \mathbf E_{\rm part} = -\,\partial_t\mathbf B,
\qquad
\nabla\times \mathbf E_{\rm hom} = 0,
\]
so that
\[
q = q_{\rm part} + q_{\rm hom}.
\]
If the homogeneous remainder is also taken to be charge-free in the bulk, then
\(\mathbf E_{\rm hom}=-\nabla\phi_{\rm hom}\) with
\(\nabla^2\phi_{\rm hom}=0\), and a harmonic DtN map becomes a natural closure
only for \(q_{\rm hom}\), not for the whole \(q\). In that reading,
\(\Lambda_U\) should be interpreted as the side operator for the homogeneous
remainder after a particular inductive contribution has been removed. The
minimal direct closure above is therefore the simplest explicit model, not the
most refined one. This refinement is legitimate in principle, but it still
does not close the problem by itself: one must choose or compute a particular
field \(\mathbf E_{\rm part}\), and hence \(q_{\rm part}\), in a way that is
consistent with the adopted boundary/interface conditions. So the split
\[
q=q_{\rm part}+q_{\rm hom}
\]
is a useful framework, not yet a unique forcing-side law.

\paragraph{Most practical experiment.}
The smallest concrete refinement of this idea is to build
\[
q = q_{\rm part} + q_{\rm hom}
\]
with the particular contribution taken from the df/inductive shell channel via
the existing nonlocal upper-side update machinery, and the homogeneous
remainder taken from the cf/curl-free shell channel via the harmonic side
operator \(\Lambda_U\). That split is now straightforward to test in the same
coefficient space as the current runtime. In the present implementation it
should still be read as experimental: with the current equivalent nonlocal
upper-side update backend, the particular contribution can be very weak on
representative pure-spectral cases, so the direct \(\Lambda_U\) closure
remains the cleaner practical baseline. A second practical variant is to keep
the same harmonic cf remainder, but take the particular contribution from the
current-first thin-sheet shell source law in \(J_S\)-space. That variant is
more closely aligned with the current runtime normalization when the forcing
branches are compared through the common shell-current input. In the present
pure-spectral implementation, that \(J_S\)-space df particular can also
collapse to zero, so this second variant is best read as a diagnostic check on
the current-first semantics rather than as a better default closure.

\paragraph{Block elimination / Schur-complement form.}
Numerically, the cleanest route is to discretize the admissible gap problem in
terms of:
\begin{itemize}
\item the shell connector coefficients \(\chi\);
\item auxiliary gap coefficients \(x\) representing whatever internal field
variables the chosen discretization uses;
\item the prescribed shell forcing coefficients \(e\) for
\(\mathbf E_{S,\rm known}\).
\end{itemize}
The condensed co-energy then has the generic quadratic form
\begin{equation}
\mathcal J(\chi,x;e)
=
\frac12
\begin{bmatrix}
\chi\\ x
\end{bmatrix}^{\!T}
\begin{bmatrix}
K_{\chi\chi} & K_{\chi x}\\
K_{x\chi} & K_{xx}
\end{bmatrix}
\begin{bmatrix}
\chi\\ x
\end{bmatrix}
-
\begin{bmatrix}
\chi\\ x
\end{bmatrix}^{\!T}
\begin{bmatrix}
M_\chi\\ M_x
\end{bmatrix}
e.
\label{eq:gamma_known_block_functional}
\end{equation}
Eliminating the internal gap variables \(x\) gives the shell-only functional
\[
\mathcal J_{\rm eff}(\chi;e)
=
\frac12 \chi^T \Lambda_{\rm gap}\chi
- \chi^T \Gamma_{\rm known} e,
\]
with the exact condensed operators
\begin{equation}
\Lambda_{\rm gap}
=
K_{\chi\chi} - K_{\chi x}K_{xx}^{-1}K_{x\chi},
\qquad
\Gamma_{\rm known}
=
M_\chi - K_{\chi x}K_{xx}^{-1}M_x.
\label{eq:gamma_known_schur}
\end{equation}
So in practical terms, computing \(\Gamma_{\rm known}\) means choosing the gap
variables \(x\), assembling \eqref{eq:gamma_known_block_functional}, and then
eliminating \(x\) by the Schur complement \eqref{eq:gamma_known_schur}. A
column-by-column solve with prescribed shell basis data is mathematically
equivalent and is the cleanest minimal implementation target:
\[
e_i \mapsto q_i,
\qquad
\mathcal D_q = [\,q_1\ \cdots\ q_N\,].
\]
Under frozen coefficients, this is just a time-independent linear response
matrix in the chosen shell basis.

The runtime-level operator interface now mirrors this abstract
Schur-complement construction explicitly: given the block data
\[
(K_{\chi\chi},K_{\chi x},K_{xx},M_\chi,M_x),
\]
it can build the condensed pair
\[
(\Lambda_{\rm gap},\Gamma_{\rm known})
\]
and then recover the corresponding \(q\)-trace by exact shell inversion. What
is still missing is not the condensation formula, but the actual PFAC/RM
assembly of those block matrices from the adopted gap model.

There is now also a smaller shell-level realization of the same interface in
the runtime: any already-condensed forcing model that supplies
\(\Gamma_{\rm known}\) can be embedded as a \emph{zero-internal-variable}
block model with
\[
K_{\chi x}=0,\qquad K_{xx}=0,\qquad M_x=0,
\]
and
\[
K_{\chi\chi}=\Lambda_{\rm shell}^{\rm exact}
=-\frac{R_I}{\mu_0}\Delta_S,\qquad
M_\chi=\Gamma_{\rm known}.
\]
This does not create a new independent gap derivation; it simply realizes an
already chosen shell forcing semantics through the same co-energy /
Schur-complement interface.

\paragraph{Why the current runtime still cannot assemble exact
\(\Gamma_{\rm known}\).}
The present PFAC/RM code already supplies important parts of the magnetic
connector:
\begin{itemize}
\item the PFAC magnetic coupling \(T\mapsto V_e\);
\item the shell-to-\(R_M\) toroidal boundary operators;
\item the exact shell-current continuity connector for the induced branch.
\end{itemize}
But that is still not enough to form \(\Gamma_{\rm known}\) exactly from the
PFAC/RM magnetic blocks \emph{alone}, because
the Schur-complement construction above requires one of the following to be
represented explicitly:
\begin{enumerate}
\item the internal gap variables \(x\) together with the quadratic co-energy
blocks \(K_{\chi\chi},K_{\chi x},K_{xx}\) and forcing blocks
\(M_\chi,M_x\); or
\item the already condensed Dirichlet-to-Neumann trace map
\[
\mathcal D_q:\mathbf E_{S,\rm known}\mapsto q_{\rm known}
\]
for \(q_{\rm known}=\partial_r U_{\rm known}-E_{r,\rm known}\).
\end{enumerate}
The PFAC/RM magnetic connector alone carries neither object. Its operators are
magnetic response blocks, not an explicit electric-trace or gap co-energy
discretization. The runtime now does expose concrete \(q\)-trace realizations
for several chosen forcing semantics, but those are additional model layers on
top of the PFAC/RM magnetic blocks, not yet the fully independent
Schur-complement assembly of \(\Gamma_{\rm known}\) from the gap model alone.
So the exact \(\Gamma_{\rm known}\) is now identified precisely, but still not
numerically assembled from the PFAC/RM magnetic blocks by themselves.

The current-first forcing branch now uses the same shell-gap realization style
as the condensed branch: the shell correction \(\Pi_{\rm shell}\) is kept as
an explicit internal gap variable, and the current-first source operator is
embedded through the same Schur-complement interface. This makes that branch a
genuine one-boundary shell-gap model rather than a purely zero-internal
condensed shell law. But the forcing-side source \(\Gamma_{\rm known}\) is
still supplied by the adopted current-first shell law
\(\delta\mathbf E_S\mapsto \delta\mathbf K_S\mapsto \delta j_r^+\), not by the
PFAC magnetic blocks alone.

\paragraph{Unified operator view.}
All of the closure routes discussed in this note can be written in the same
shell-source form:
\begin{equation}
-\,\frac{R_I}{\mu_0}\Delta_S\chi
=
s_{\rm known} + s_{\rm ind},
\qquad
\chi:=\partial_t\psi^+.
\label{eq:unified_shell_source_view}
\end{equation}
What differs between the routes is not the final shell inversion, but how the
source terms are built. The induced branch uses the thin-sheet continuity
connector, the default condensed forcing branch uses the reduced inductive
source operator, the frozen-conductance branch uses the current-first
incremental shell law
\[
\delta\mathbf K_S \mapsto \delta j_r^+
\]
with a frozen-conductance adapter
\[
\delta\mathbf E_S \mapsto \delta\mathbf K_S,
\]
and the abstract gap-BVP/co-energy route would define the same source through
the condensed operator \(\Gamma_{\rm known}\). The appearance of
\(\partial_t\)-notation in some places and \(\delta\)-notation in others
reflects the same distinction: \(\partial_t\) is the continuous-time shell
equation, while \(\delta\) denotes the linearized connector-substep form of
the same source-building problem.

Across all of these routes, the common shell-facing returned quantity is
\(\chi\), with \(q\) and \(s^+\) serving only as its exact forcing-side
preimage and source-side image:
\[
q \xmapsto{\ \chi=(1/R_I)\Pi_0 q\ } \chi
\xmapsto{\ s^+=-(R_I/\mu_0)\Delta_S\chi\ } s^+.
\]
The distinction between the routes lies only in how they build that returned
shell variable or its exact images.

For the finite-\(R_M\) gap picture, it is therefore natural to separate the
outer magnetic data from the ionospheric connector solve:
\[
q = \mathcal D_q^{\rm iono}\,\mathbf E_{S,\rm shell}
  + \mathcal D_q^{RM}
  \begin{bmatrix}
    \psi^+(R_M)\\
    M(R_M)
  \end{bmatrix}.
\]
The first term is the pure ionosphere-to-connector map, while the second term
represents prescribed magnetospheric boundary forcing through the finite-\(R_M\)
magnetic pair.

The smallest constructive implementation target is therefore not the full
combined \(4N\to N\) map at once, but the homogeneous-outer-data operator
\[
\mathcal D_q^{\rm iono}:
\begin{bmatrix}
U_I\\
V_I
\end{bmatrix}
\mapsto q_I
\]
with
\[
\psi^+(R_M)=0,
\qquad
M(R_M)=0.
\]
That operator can be assembled column by column. In the current runtime, the
first implemented version of that construction uses the existing harmonic
side-trace closure as the per-column baseline, so it is best read as a
constructive shell-level realization of the homogeneous finite-\(R_M\)
closure target, not yet as a fully independent gap discretization.

Operationally, both the forcing and induced branches can now be viewed in the
same intermediate variables:
\[
\text{input}
\;\to\;
\partial_t j_r^+
\;\to\;
q=R_I\,\partial_t\psi^+
\;\to\;
\chi=\partial_t\psi^+.
\]
What changes from one branch to another is only the source-building map from
the chosen input data to \(\partial_t j_r^+\).

So the most uniform shell-level abstraction is not to force every branch to
share the same \emph{input} variable, but to require every branch to expose
the same condensed operator bundle:
\[
\Lambda_{\rm gap},
\qquad
\Gamma_{\rm known},
\qquad
\mathcal D_q,
\]
with optional current-first and induced-response counterparts when those are
available. In that view, the solved-for quantity is always
\(\chi=\partial_t\psi^+\), the exact shell operator is always
\(\Lambda_{\rm shell}^{\rm exact}=-(R_I/\mu_0)\Delta_S\), and the only
branch-dependent part is which source-builder produces
\(\Gamma_{\rm known}\) or \(\partial_t j_r^+\).

For the \emph{current runtime forcing branches}, one can now go one step
further and normalize the interface to the same practical shell input
variable:
\[
\mathbf J_S \mapsto \Gamma_{\rm known}^{(J)} \mapsto q \mapsto \chi.
\]
Branches that are naturally current-first, such as the frozen-conductance
path, expose that map directly. Branches that are naturally electric-first,
such as the reduced condensed-inductive path, now expose the same input by
composition with the live shell constitutive operator
\[
\mathbf J_S \mapsto \mathbf E_S.
\]
So the present forcing-side runtime can be read uniformly as
\emph{shell-current input} \(\to\) source \(\to q \to \chi\), even though the
underlying source-building semantics still differ. This should be read as an
\emph{interface-level normalization}, not as a claim that every branch is
natively current-driven in the same physical sense.

Because of the zero-internal-variable realization just described, the current
runtime forcing branches can now also be viewed as shell-level co-energy
realizations of their own condensed source operators. That unifies the
operator interface across the branches, even though only a future explicit
gap-block assembly would count as a genuinely independent PFAC/RM
construction.

\paragraph{Where energy minimization fits.}
The co-energy or energy-minimization viewpoint is therefore not a missing extra
equation for the frozen-conductance branch. It is an \emph{alternative source
builder} for the same forcing-side quantity \(s_{\rm known}\). In the present
notation:
\begin{itemize}
\item the condensed inductive runtime branch builds \(s_{\rm known}\) from the
reduced projected-tangential operator;
\item the frozen-conductance branch builds \(s_{\rm known}\) from a
current-first shell law plus a constitutive adapter;
\item the gap-BVP/co-energy route would build \(s_{\rm known}\) by solving the
adopted upper-region problem and condensing it to \(\Gamma_{\rm known}\).
\end{itemize}
So co-energy is best understood as a more ambitious \emph{alternative source
builder}, not as a correction factor that the shell-current branch is still
missing internally.

\paragraph{Notation convention.}
The cleanest convention is therefore:
\begin{itemize}
\item use \(\partial_t\) for the exact shell identities, the gap BVP, and the
shell-condensed closure operators;
\item use \(\delta\) only for reduced incremental laws such as the
frozen-conductance forcing branch, where the model is justified specifically as
a connector-substep increment law rather than as a general exact
continuous-time identity;
\item keep variational perturbations \(\delta\chi\) local to the variational
paragraphs, where they denote test variations of the shell functional rather
than physical time increments.
\end{itemize}
If one freezes coefficients over a short substep \(\Delta\tau\), then the two
physical views are related by
\[
\delta X \approx (\partial_t X)\,\Delta\tau,
\]
but they should not be conflated in the theory text because the frozen
conductance branch is only claimed as an incremental closure.

\section*{Conclusion}

\paragraph{Exact statement.}
The canonical scalar closure is the radial-shell equation. In its cleanest
form, the missing upper-side quantity is one scalar connector variable,
equivalently
\[
\partial_t\psi^+,
\qquad
\partial_t j_r^+,
\qquad
Q_I-\Delta_\Omega E_{r,I}.
\]
Once that scalar is known, the final shell inversion is exact.

\paragraph{Model statement.}
The note now separates the \emph{reduced-shell} closure object from the
\emph{forcing-side} builder used to reach it. At reduced-shell level, the
primary nonlocal object is
\[
\Pi = \mathcal P_I[\chi],
\]
and the closure is the boundary equation
\[
K_I \chi = -R_I \widetilde G.
\]
When the primitive runtime input is the shell electric forcing rather than
\(\chi\) itself, the same upper response is organized instead through the
forcing-side source builder \(\Gamma_{\rm known}\), or equivalently the shell
trace map \(\mathcal D_q\) that first returns \(q\) and then maps exactly to
\(\chi\). The note now makes clear which versions of these objects are:
\begin{itemize}
\item exact shell consequences;
\item explicit model choices, such as harmonic trace continuation or the
frozen-conductance incremental forcing law;
\item reduced runtime realizations, such as the default condensed inductive
forcing path.
\end{itemize}
So the remaining gap is now narrow and concrete: the constructive assembly of
the condensed upper-region operator for the chosen boundary semantics,
discretization, and shell/gap interface data, whether one writes that
operator as \(\mathcal P_I\) in the reduced-shell magnetic language or as
\(\Gamma_{\rm known}/\mathcal D_q\) in the forcing-side electric language.

\paragraph{Runtime statement.}
The current code has an explicit PFAC/RM radial-shell response path, not only
the older projected tangential route. In that runtime:
\begin{itemize}
\item the feedback side is closed by the explicit thin-shell
current-continuity connector;
\item the forcing side is exposed through the exact shell quantity
\[
q_{\rm known}=\partial_r U_{\rm known}-E_{r,\rm known};
\]
\item the magnetic runtime now distinguishes the raw PFAC return
\(\Pi_{\rm raw}\), the harmonic shell baseline determined by
\(B_r(R_I)\), and the baseline-subtracted shell correction
\(\Pi_{\rm shell}\) with zero shell radial trace;
\item the present forcing-side runtime can be read uniformly as a
shell-current-input path
\[
\mathbf J_S \to \partial_t j_r^+ \to \chi,
\]
with \(q\) retained only as the exact intermediate preimage when one wants the
electric trace explicitly, and with electric-first branches exposed through the
live constitutive adapter \(\mathbf J_S \mapsto \mathbf E_S\);
\item the built-in runtime now keeps a single supported forcing-side
realization, \path{frozen_conductance_incremental}, while older condensed and
tangential benchmark routes are retained only as internal diagnostics.
\end{itemize}

\paragraph{Final status.}
The strongest honest summary is therefore: the exact shell backbone is
established; the primary reduced-shell closure object is
\[
\Pi=\mathcal P_I[\chi]
\]
with \(\chi=\partial_t\psi^+\); the shell-facing returned quantity is
\(\chi=\partial_t\psi^+\) or equivalently \(s^+=\partial_t j_r^+\); the
forcing-side maps \(\Gamma_{\rm known}\) and \(\mathcal D_q\) are the
constructive interfaces that build the same connector from shell forcing data;
the independent completion of the shell model is still a condensed
upper-region BVP or parent transmission solve; and what remains is the
constructive assembly of that operator, with precise input semantics fixed for
the chosen runtime or discretization. There is no hidden local shell law that
closes full induction by itself.

\appendix

\section{The current built-in PFAC/RM radial-shell model}

The current code now contains an explicit built-in
\path{radial_shell} runtime path. It is not a future placeholder anymore.
Its structure is:
\begin{equation}
\eta
\;\to\;
\chi
\;\to\;
[\Phi, W]
\;\to\;
\text{radial-shell scalar response},
\label{eq:runtime_radial_chain}
\end{equation}
with a forcing-side branch, a feedback-side branch, and optional explicit
trace-based variants. Internally the current code may still pass through
equivalent connector or driver variables, but the reduced-shell reading of the
runtime should now be understood as \(\chi\)-first.

\paragraph{How this appendix should be read.}
The appendix is still organized around the runtime scalar connector path,
because that is how the present code is assembled operationally. But after the
updates described above, the same runtime should also be read in the magnetic
language of the companion document, with \(\chi\) as the primary shell
variable and \(q\) only as its forcing-side preimage:
\[
\text{shell forcing} \to \chi
\qquad
(\text{with } q = R_I \chi \text{ as the forcing-side preimage})
\qquad
\to \Pi_{\rm raw} \to \Pi_{\rm harm} \to \Pi_{\rm shell}.
\]
So the runtime appendix is no longer merely a discussion of \(\Gamma_{\rm known}\)
and \(\mathcal D_q\); it is also the place where the code-level realization of
the reduced-shell operator \(\Pi=\mathcal P_I[\chi]\) is identified.

\subsection*{1. Optional shell-electric trace model}

The runtime stores shell tangential electric coefficients in the active
Helmholtz basis,
\[
[\Phi, W].
\]
One explicit route to the forcing-side radial-shell scalar is the trace model
\[
[\Phi, W]
\mapsto
\left(
\partial_r U|_{R_I^+},
E_{r,I}
\right)
\mapsto
-\frac{1}{R_I^2}\Bigl[Q_I - \Delta_\Omega E_{r,I}\Bigr].
\]
This optional trace route makes the following explicit assumptions:
\begin{enumerate}
\item the physical curl-free shell electric scalar is
\[
U = \mathrm{cf\_sign}\,\Phi;
\]
\item only that poloidal electric scalar is radially continued to
$\partial_r U|_{R_I^+}$;
\item the toroidal electric shell scalar $W$ does \emph{not} enter the
$\partial_r U$ continuation;
\item the shell-normal electric trace is recovered from the full shell
tangential field by the ideal-gap relation
\[
E_r = -\,\frac{\mathbf B_{0S}\cdot \mathbf E_S}{B_{0r}},
\qquad
(B_{0r}\neq 0),
\]
with low-latitude/small-\(B_{0r}\) treatment handled through the source or
connector formulation rather than by pointwise division.
\end{enumerate}

\paragraph{Relation to PFAC.}
The shell electric coefficients themselves are not local objects. They are
already obtained from the live conductivity/PFAC pathway, including the dynamic
\texttt{psi -> Ve} response and the chosen $R_M$ open/closed shielding
semantics. So the trace model is applied to the full shell electric response
already produced by the code, not to a shell-local proxy field.

\paragraph{Incremental ideality view.}
There is a useful equivalent way to think about this trace route. Since the
background field is static, ideality implies
\[
\partial_t(\mathbf E\cdot \mathbf B_0)=0
\quad\Longrightarrow\quad
(\partial_t\mathbf E)\cdot \mathbf B_0 = 0,
\]
and therefore, where $B_{0r}\neq 0$,
\[
\partial_t E_r
=
-\,\frac{\mathbf B_{0S}\cdot \partial_t\mathbf E_S}{B_{0r}}.
\]
Combining this with the toroidal part of Faraday's law gives the increment form
\[
R_I\,\delta\psi^+ = \Pi_0(\partial_r \delta U - \delta E_r),
\]
so once a consistent trial pair
\[
\bigl(\delta \mathbf E_S,\ \delta\psi^+\bigr)
\]
is known, one may recover the increment in the missing radial-shell scalar
without introducing a separate raw continuation field:
\[
\delta Q_I
=
R_I\,\Delta_\Omega \delta\psi^+ + \Delta_\Omega \delta E_r.
\]
This is a useful \emph{correction/increment} viewpoint on the same linear
closure. It does not provide a new standalone closure law by itself; it still
depends on having a consistent upper-side connector response for the same trial
update.

\subsection*{2. Built-in forcing-side q-trace models}

The current built-in \path{radial_shell} runtime now routes the forcing side
through the exact shell quantity
\[
q_{\rm known} = \partial_r U_{\rm known} - E_{r,\rm known}.
\]
What varies with the selected forcing semantics is only \emph{how} that
\(q\)-trace is constructed.

\paragraph{Optional direct side-trace path.}
There is now also a smaller optional explicit route for the minimal direct
closure from Section 10:
\[
U|_{S_I}
\xmapsto{\ \Lambda_U\ }
\partial_r U|_{R_I^+},
\qquad
q_{\rm known}=\Lambda_U U-E_{r,\rm known}.
\]
In the current code, open harmonic and \(R_M\)-shielded harmonic realizations
of \(\Lambda_U\) are available as explicit direct side operators. This is an
optional explicit trace route, not the built-in default forcing semantics.
At present that route models the whole \(q\) directly; a more refined future
variant could subtract a particular inductive contribution first and apply the
harmonic side operator only to the homogeneous remainder.

\paragraph{Default condensed-inductive path.}
The default built-in forcing mode still uses the already assembled nonlocal
PFAC/projected-tangential full-induction response, but it is now expressed in
the exact one-curl shell language:
\[
\mathbf E_{\rm known}
\mapsto
\partial_t\psi_{\rm known}^+
\mapsto
q_{\rm known}=R_I\,\partial_t\psi_{\rm known}^+
\mapsto
\partial_t j_{r,\rm known}^+.
\]
Operationally, the same reduced tangential solve still sits underneath that
update, so the concrete chain is
\begin{equation}
\begin{aligned}
\mathbf E_{\rm known}
&\xmapsto{\ f_{t,\mathrm{red}}\ }
\text{projected tangential RHS}
\xmapsto{\ (L_\alpha^{\rm tang})^{-1}\ }
\partial_t\alpha_{\rm tangential}
\\
&\xmapsto{\ \alpha\mapsto \partial_t\psi_{\rm known}^+\ }
q_{\rm known}=R_I\,\partial_t\psi_{\rm known}^+
\xmapsto{\ -(\mu_0 R_I^2)^{-1}\Delta_\Omega\ }
\partial_t j_{r,\rm known}^+.
\end{aligned}
\label{eq:built_in_condensed_forcing_chain}
\end{equation}
So the default forcing path is still a reduced-model realization of the exact
radial-shell forcing source, but its final shell assembly is now written in
the exact \(q\)-trace language instead of only as
\(\mu_0 B_{0r}\eta\). In the current code this branch is also routed through
the same Schur-complement adapter as a degenerate shell-level co-energy model,
so the runtime now speaks uniformly in
\[
\Gamma_{\rm known},\qquad q_{\rm known},\qquad
\Lambda_{\rm shell}^{\rm exact}
\]
even though the forcing semantics underneath remain the same reduced
projected-tangential ones as before.

\paragraph{Alternative frozen-conductance path.}
The alternative built-in forcing mode
\path{frozen_conductance_incremental} now passes through the same exact
\(q\)-trace interface, but by a different realization:
\[
\delta \mathbf E_{S,\rm known}
\mapsto
\delta \mathbf K_{S,\rm known}
\mapsto
\delta j_{r,\rm known}^+
\mapsto
\delta\psi_{\rm known}^+
\mapsto
q_{\rm known}=R_I\,\delta\psi_{\rm known}^+.
\]
Here the primitive source law is the current-first thin-sheet continuity
relation
\[
\delta \mathbf K_{S,\rm known}\mapsto \delta j_{r,\rm known}^+,
\]
while the frozen-conductance electric step
\[
\delta \mathbf E_{S,\rm known}\mapsto \delta \mathbf K_{S,\rm known}
\]
is only the constitutive adapter placed on top. The \(q\)-trace is then
obtained by exact shell inversion rather than by the reduced tangential update
route.

\paragraph{Source-level view.}
In the newer shell-source language, both supported forcing modes are now just
two realizations of the same exact interface
\[
\mathcal D_q:\mathbf E_{S,\rm known}\mapsto q_{\rm known},
\]
and therefore of the same shell-current source operator
\begin{equation}
\mathbf E_{\rm known}
\mapsto
s_{\rm known}^+
:=
\partial_t j_{r,\rm known}^+
=
\frac{1}{\mu_0}\,\mathcal R_{\rm known}\,\mathbf E_{\rm known},
\label{eq:built_in_known_source_operator}
\end{equation}
with the built-in \path{frozen_conductance_incremental} path using exact
source inversion. What is still not available from the PFAC/RM magnetic connector
alone is a fully independent constructive derivation of \(\mathcal D_q\) from
the gap model without going through one of those chosen forcing realizations.

\subsection*{3. Feedback-side upper toroidal update model}

For the induced feedback, the runtime now works directly with the missing
scalar rather than only an equivalent condensed tangential adapter. In the
built-in path the decisive identity is
\[
\mu_0\,\partial_t j_r^+
=
-\,R_I \Delta_S(\partial_t\psi^+),
\]
so the feedback branch is best described at the source level rather than as a
separate toroidal-surrogate solve.

\paragraph{Source-level induced operator.}
With the present continuity-based connector, the induced upper-side source is
now explicit:
\begin{equation}
\partial_t j_{r,\rm ind}^+
=
\partial_t j_r - \nabla_S\cdot \partial_t\mathbf K_S.
\label{eq:continuity_induced_source_shell}
\end{equation}
In coefficient form this is
\begin{equation}
\mathcal S_{\rm ind}
=
T_{\alpha\to j_r}
-\frac{1}{R_I}\,\mathrm{div}_\Omega\,
T_{\psi\to J_S}\,
T_{\alpha\to\psi},
\label{eq:continuity_induced_source_operator}
\end{equation}
so
\[
\partial_t j_{r,\rm ind}^+
=
\mathcal S_{\rm ind}\,\partial_t\alpha.
\]
This is the precise PFAC/RM shell-source operator for the induced branch in
the current runtime.

\paragraph{Combined shell-source closure.}
The built-in radial-shell mode can therefore be written directly in source
form as
\begin{equation}
-\,\frac{R_I}{\mu_0}\Delta_S(\partial_t\psi^+)
=
s_{\rm known}^+(\mathbf E_{\rm known})
+
\mathcal S_{\rm ind}\,\partial_t\alpha.
\label{eq:runtime_shell_source_split}
\end{equation}
In this form, the exact final inversion is the shell Laplacian, while the
model dependence sits entirely in the known forcing source
\(s_{\rm known}^+\) and, if one changed connector models, in the source
builder \(\mathcal S_{\rm ind}\).

\subsection*{4. Automatic PFAC/RM mode selection}

The built-in \path{radial_shell} runtime chooses the explicit open/closed
response branch from the same runtime assumptions used elsewhere:
\begin{itemize}
\item when $R_M$ is absent, the shell trace and upper toroidal update both use
the open branch;
\item when $R_M$ is present and magnetospheric shielding is enabled, the shell
trace uses the shielded harmonic continuation and the toroidal update uses the
closed $R_M$ branch.
\end{itemize}
So the default runtime \path{radial_shell} mode is now an explicit PFAC/RM
radial-shell model, not a generic placeholder.

\subsection*{5. Candidate explicit RM/PFAC connector system}

The existing code already contains most of a candidate explicit connector
system. Written as a chain in coefficient space, the physically unambiguous
parts are:
\begin{align}
\eta
&\mapsto
\partial_t j_r(R_I)
=
B_{0r}(R_I)\,\eta,
\\
\eta
&\mapsto
\partial_t\psi(R_I)
=
T_{\alpha\to\psi}\,\eta,
\\
\eta
&\mapsto
j_n(R_M)
=
T_{\alpha\to j_n^{RM}}\,\eta,
\\
\Delta_S \chi(R_M)
&=
j_n(R_M),
\\
\partial_t\psi_b(R_M)
&=
\frac{\mu_0}{R_M}\,\chi(R_M),
\\
\partial_t\psi_{\rm raw}^+(R_I)
&=
T_{\alpha\to\psi}(R_I)\,\eta
\;+\;
S_{RM\to RI}^{\rm int}\,\partial_t\psi_b(R_M),
\\
\partial_t V_e(R_I)
&=
T_{\psi\to V_e}^{\rm PFAC}\,\partial_t\psi(R_I).
\end{align}

All of these operators already exist explicitly in the runtime:
\begin{itemize}
\item \(T_{\alpha\to j_n^{RM}}\), \(\chi(R_M)\), and
\(\partial_t\psi_b(R_M)\) are the toroidal \(R_M\) closure operators;
\item \(S_{RM\to RI}^{\rm int}\) is the internal radial continuation from
\(R_M\) back to \(R_I\);
\item \(T_{\psi\to V_e}^{\rm PFAC}\) is the dynamic PFAC operator mapping the
toroidal channel to the nonlocal poloidal magnetic response.
\end{itemize}

\paragraph{How this now fits the document's \(\Pi\) operator.}
The runtime now distinguishes two related objects:
\[
\Pi_{\rm raw}^{\rm runtime}:=T_{\psi\to V_e}^{\rm PFAC},
\]
the toroidal-to-poloidal magnetic return selected with the live open/closed
\(R_M\) semantics, and
\[
\Pi_{\rm shell}^{\rm runtime}
:=
\Pi_{\rm raw}^{\rm runtime}-\Pi_{\rm harm}^{\rm runtime},
\qquad
B_r\bigl(\Pi_{\rm shell}^{\rm runtime}\bigr)\big|_{R_I}=0,
\]
the document-style shell correction after subtraction of the harmonic upward
baseline fixed by the returned shell radial trace. So the raw PFAC block is
the magnetic \(\Pi\)-ingredient, while the reduced-shell operator of the
companion document is more naturally identified with the baseline-subtracted
shell correction. The code also exposes the realization of the raw block
through the full shell chain
\[
\Pi_{\rm raw}^{\rm runtime}
\;\to\;
\bigl(\partial_t\psi \mapsto J_S\bigr)
\;\to\;
\bigl(\partial_t\psi \mapsto E_S\bigr)
\;\to\;
\bigl(\partial_t\psi \mapsto \text{toroidal rhs}\bigr)
\;\to\;
\bigl(\partial_t\psi \mapsto \dot\psi\bigr).
\]

\paragraph{Open vs.\ shielding split.}
For both objects the runtime makes the same open-plus-shielding split explicit:
\[
\Pi_{{\rm raw},\,\rm eff}^{\rm runtime}
=
\Pi_{{\rm raw},\,\rm open}^{\rm runtime}
+
\Pi_{{\rm raw},\,\rm shield}^{\rm runtime},
\]
where \(\Pi_{{\rm raw},\,\rm shield}^{\rm runtime}\) is exactly the
reflected/internal \(R_M\)-closure reaction used for magnetospheric shielding.
At the shell itself, the code now also exposes the corresponding
baseline-subtracted correction split
\[
\Pi_{{\rm shell},\,\rm eff}^{\rm runtime}
=
\Pi_{{\rm shell},\,\rm open}^{\rm runtime}
+
\Pi_{{\rm shell},\,\rm shield}^{\rm runtime},
\qquad
B_r\bigl(\Pi_{\rm shell}^{\rm runtime}\bigr)\big|_{R_I}=0.
\]
Here \(\Pi_{{\rm raw},\,\rm shield}^{\rm runtime}\) is exactly the
reflected/internal \(R_M\)-closure reaction used for magnetospheric shielding.

\paragraph{From raw \(\Pi\) to the shell correction.}
In the present implementation, \(V_e\) is a poloidal magnetic potential, not
the radial field itself. So the exact next layer is
\[
\Pi_{\rm raw}^{\rm runtime}
\;\to\;
V_e
\;\to\;
B_r,
\]
using the same magnetic scalar-to-\(B_r\) map that the code already uses for
the induced poloidal potential \(m_{\rm ind}\). Once \(B_r(R_I)\) is known,
the harmonic upward baseline is the unique shell scalar with the same radial
trace there; subtracting it yields \(\Pi_{\rm shell}^{\rm runtime}\). The
runtime now exposes this whole chain, including the open/effective/shielding
splits at both \(R_I\) and \(R_M\).

\paragraph{Radial derivative at \(R_I^+\).}
The runtime now also exposes the corresponding
\[
\Pi_{\rm raw}^{\rm runtime}
\;\to\;
\partial_r B_r(R_I^+)
\]
operator. This is assembled mode by mode from the same harmonic branch split:
the open branch is treated as the external contribution at \(R_I^+\), while
the shielding branch is treated as the reflected/internal contribution there.
So at the current implementation level, \(\partial_r B_r(R_I^+)\) is now an
explicit operator derived from the open-plus-shielding decomposition, not a
numerical derivative of the PFAC radial integrals.

At \(R_M\) the branch identification is different: the open PFAC return there
is internal relative to the evaluation radius, because the gap sources lie
below the truncation boundary. The runtime now exposes the corresponding
\[
\Pi_{\rm raw}^{\rm runtime}
\;\to\;
\partial_r B_r(R_M^-)
\]
operator with that internal-branch convention built in, again as an explicit
operator split rather than a differentiated radial quadrature.

\paragraph{Arbitrary-radius split.}
The same construction is now available for any chosen
\[
R_I \le r \le R_M.
\]
Relative to that evaluation radius, the open PFAC return is split into the
direct internal and external pieces coming from source layers below and above
\(r\), respectively. The effective shielded operator is then assembled from
that direct split plus the reflected \(R_M\)-closure contribution, using the
same harmonic roundtrip factor already present in the PFAC kernel. So the
runtime now exposes a genuinely radius-dependent \emph{raw} operator family
\[
\Pi_{\rm raw}(r),\qquad B_r(r),\qquad \partial_r B_r(r),
\]
rather than only the two endpoint special cases \(R_I^+\) and \(R_M^-\).
The baseline-subtracted shell correction \(\Pi_{\rm shell}\) is naturally an
\(R_I\)-anchored object, because its defining harmonic subtraction is fixed by
the shell radial trace \(B_r(R_I)\).

 \paragraph{What is explicit already, and what is still missing.}
The explicit RM/PFAC chain already gives a genuine magnetic connector picture:
a toroidal/FAC branch through the incoming normal current and its \(R_M\)
closure current, a raw upper-side toroidal update
\(\partial_t\psi_{\rm raw}^+\), a coupled PFAC/poloidal magnetic response
\(\partial_t V_e\), and the shell harmonic baseline together with the
baseline-subtracted correction \(\Pi_{\rm shell}^{\rm runtime}\) at \(R_I\).
What it still does \emph{not} give by itself is a closed independent
derivation of the full reduced-shell operator \(\mathcal P_I\), because it
does not yet solve an independent frozen gap transmission/BVP and then
condense it. The missing step is the constitutive completion that turns the
magnetic pair
\[
\bigl(\partial_t\psi_{\rm raw}^+,\ \partial_t V_e\bigr)
\;\to\;
\partial_t \mathbf K_S
\;\to\;
\partial_t \mathbf E_S
\;\to\;
\partial_t\psi^+,
\]
that is, magnetic connector response completed by shell jump/current balance,
frozen shell electrodynamics, ideality, and the toroidal part of Faraday's
law. The current runtime already exposes the shell-electrodynamic part of this
chain explicitly as
\[
\partial_t\psi \to \partial_t\mathbf J_S \to \partial_t\mathbf E_S,
\]
but the fully independent upper-region completion law is still not assembled.

\paragraph{What is still condensed.}
The feedback connector itself need not remain condensed. If the ionosphere is
treated as an instantaneous thin current sheet with no stored surface charge,
then sheet-current continuity gives
\[
\nabla_S\cdot \partial_t\mathbf K_S + \partial_t j_r^+ - \partial_t j_r = 0.
\]
Since \(\partial_t j_r = B_{0r}\eta\), the missing upper-side scalar is fixed
by the explicit connector law
\begin{equation}
\partial_t j_r^+
=
B_{0r}\eta - \nabla_S\cdot \partial_t\mathbf K_S.
\label{eq:thin_shell_current_continuity_connector}
\end{equation}
Equivalently, in coefficient form with the angular divergence operator,
\[
\partial_t j_r^+
=
B_{0r}\eta - \frac{1}{R_I}\,\mathrm{div}_\Omega(\partial_t\mathbf K_S).
\]
This closes the feedback connector directly once
\(\partial_t\psi\mapsto\partial_t\mathbf K_S\) is known.

\paragraph{Why this part is robust.}
Within the adopted instantaneous thin-sheet/no-stored-charge shell model,
\eqref{eq:thin_shell_current_continuity_connector} is simply the pillbox
current-balance identity for the collapsed ionospheric sheet. Over the
instantaneous connector substep, radial current can only:
\begin{enumerate}
\item arrive from below as the shell-side FAC increment \(B_{0r}\eta\);
\item spread laterally inside the sheet as \(\partial_t\mathbf K_S\); or
\item leave upward as the unknown connector current \(\partial_t j_r^+\).
\end{enumerate}
So \eqref{eq:thin_shell_current_continuity_connector} simply says
\[
\text{upward outgoing current}
=
\text{incoming lower FAC}
-
\text{current diverted sideways in the sheet}.
\]
That is the clearest physical interpretation of the feedback connector, and it
is why thin-sheet current continuity belongs to the strong shell-model
backbone rather than to the fragile forcing-side modeling choices.

\paragraph{Exact shell variational principle for the connector scalar.}
Once the outgoing upper-side radial-current source
\[
s := \partial_t j_r^+
=
B_{0r}\eta - \nabla_S\cdot \partial_t\mathbf K_S
\]
is known, the upper-side toroidal update
\[
\chi := \partial_t\psi^+
\]
is equivalently the unique mean-zero minimizer of the strictly convex shell
functional
\begin{equation}
\mathcal J_{\rm shell}[\chi]
:=
\frac{R_I}{2\mu_0}
\int_{S_I} |\nabla_S \chi|^2\,dS
-\int_{S_I} s\,\chi\,dS,
\qquad
\int_{S_I}\chi\,dS = 0.
\label{eq:shell_variational_connector}
\end{equation}
Varying \(\chi\to \chi+\varepsilon\,\delta\chi\) with mean-zero
\(\delta\chi\) gives
\[
\delta \mathcal J_{\rm shell}
=
\int_{S_I}
\left[
-\,\frac{R_I}{\mu_0}\Delta_S\chi - s
\right]\delta\chi\,dS.
\]
So the Euler--Lagrange condition is
\begin{equation}
-\,\frac{R_I}{\mu_0}\Delta_S(\partial_t\psi^+)
=
\partial_t j_r^+,
\label{eq:shell_variational_connector_EL}
\end{equation}
which is exactly the connector identity
\[
\mu_0\,\partial_t j_r^+ = -\,R_I\Delta_S(\partial_t\psi^+).
\]
This is the cleanest shell-only variational statement of the completed
connector. It does \emph{not} determine the source \(s\) by itself; rather, it
packages the exact final inversion from the upper-side radial-current source to
the toroidal connector scalar.

\paragraph{What this variational form does and does not buy us.}
The variational principle \eqref{eq:shell_variational_connector} is exact once
the source \(s\) is specified. It therefore reformulates the connector cleanly,
but it is not a new independent forcing-side closure law. What remains
condensed is the forcing-side scalar response from known shell electric
forcing. The code is now explicit through
\[
\partial_t\psi \to \partial_t\mathbf J_S \to \partial_t\mathbf E_S,
\]
\[
\eta \to \partial_t j_r^+ \to \partial_t\psi^+,
\]
and the only remaining reduced-model step is the map from known shell electric
forcing to the radial-shell scalar residual. That forcing-side map is still
represented by the chosen nonlocal scalar-response model rather than by an
independent gap-region PDE derivation.

\paragraph{What the current condensed forcing operator actually is.}
The surviving nontrivial forcing object in the runtime is not a shell-current
continuity law. It is the reduced inductive tangential forcing scalar used by
the current projected-shell toroidal closure. If
\[
\mathbf E_S = E_\theta\,\hat\theta + E_\phi\,\hat\phi,
\]
define the surface-curl scalar
\begin{equation}
c
:=
\mathrm{curl}_\Omega(\mathbf E_S)
=
\partial_\theta E_\phi + \cot\theta\,E_\phi - \partial_\phi^\ast E_\theta.
\label{eq:condensed_forcing_surface_curl}
\end{equation}
Then the live condensed forcing operator is
\begin{equation}
f_{\rm ind}(\mathbf E_S)
=
\frac{1}{R_I^2}
\left[
B_{0\theta}\,\partial_\phi^\ast c
- B_{0\phi}\,\partial_\theta c
\right].
\label{eq:condensed_forcing_inductive_scalar}
\end{equation}
This is exactly the forcing block assembled by the code before coefficient
projection.

\paragraph{Interpretation.}
Let
\[
\omega_E := \hat r\cdot(\nabla\times \mathbf E)
= \frac{1}{R_I}\,c
\]
denote the radial curl scalar of the shell tangential field. Then
\eqref{eq:condensed_forcing_inductive_scalar} is equivalently
\begin{equation}
f_{\rm ind}(\mathbf E_S)
=
\mathbf B_{0S}\cdot
\nabla\times\!\bigl(\omega_E\,\hat r\bigr)\Big|_{r=R_I}.
\label{eq:condensed_forcing_toroidal_vector_form}
\end{equation}
So the live forcing keeps only the tangential toroidal piece generated by the
radial curl scalar \(\omega_E\). It discards the remaining tangential pieces of
\(-\nabla\times\nabla\times \mathbf E\), including explicit radial-derivative
information. Equation \eqref{eq:condensed_forcing_inductive_scalar} therefore
belongs to the reduced tangential closure, not to the exact radial-shell
equation.

\paragraph{Exact split of the projected Maxwell driver.}
Write
\[
\mathbf F := \nabla\times \mathbf E
=
\omega_E\,\hat r + \mathbf F_S.
\]
Then the exact projected tangential Maxwell driver splits as
\begin{equation}
f_t^{\rm drv}
=
\mathbf B_{0S}\cdot\bigl[-\nabla\times \mathbf F\bigr]_S
=
\mathbf B_{0S}\cdot\bigl[-\nabla\times(\omega_E\hat r)\bigr]_S
\;+\;
\mathbf B_{0S}\cdot\bigl[-\nabla\times \mathbf F_S\bigr]_S.
\label{eq:projected_driver_split}
\end{equation}
The first term in \eqref{eq:projected_driver_split} is exactly the live
condensed forcing block \eqref{eq:condensed_forcing_toroidal_vector_form}. The
second term,
\begin{equation}
f_{t,\rm omitted}^{\rm drv}
:=
\mathbf B_{0S}\cdot\bigl[-\nabla\times \mathbf F_S\bigr]_S,
\label{eq:projected_driver_omitted_term}
\end{equation}
is the tangential-curlcurl content dropped by the reduced runtime forcing
operator.

\paragraph{What lives in the omitted term.}
Using the shell electric decomposition
\[
\mathbf E
=
E_r\,\hat r + \nabla_S U - \hat r\times \nabla_S V,
\]
the tangential part of \(\nabla\times \mathbf E\) is
\[
\mathbf F_S
=
\hat r\times \nabla_S(\partial_r U - E_r)
+
\nabla_S(\partial_r V).
\]
So the omitted term \eqref{eq:projected_driver_omitted_term} contains exactly
the part of the projected Maxwell driver that depends on the off-shell electric
trace information \(\partial_r U - E_r\) and \(\partial_r V\). This makes the
status of the built-in forcing route precise:
\begin{itemize}
\item the live operator keeps the radial-curl/toroidal part of the exact
projected driver;
\item it omits the part requiring explicit off-shell electric-trace
information.
\end{itemize}

More explicitly, if
\[
q := \partial_r U - E_r,
\qquad
p := \partial_r V,
\]
then the omitted projected driver is
\begin{equation}
f_{t,\rm omitted}^{\rm drv}
=
\mathbf B_{0S}\cdot \nabla_S(\partial_r q)
+
\mathbf B_{0S}^{\perp}\cdot \nabla_S(\partial_r p),
\label{eq:projected_driver_omitted_second_trace_form}
\end{equation}
or equivalently
\begin{equation}
f_{t,\rm omitted}^{\rm drv}
=
\mathbf B_{0S}\cdot \nabla_S\!\bigl(\partial_r(\partial_r U - E_r)\bigr)
+
\mathbf B_{0S}^{\perp}\cdot \nabla_S(\partial_r^2 V).
\label{eq:projected_driver_omitted_second_trace_expanded}
\end{equation}
So the omitted forcing contribution depends on \emph{second radial trace
information}, not only on the first shell traces used by the optional
\((\partial_r U,E_r)\) models. This is why the present trace machinery is not,
by itself, enough to reconstruct the full projected driver.

Equation \eqref{eq:condensed_forcing_inductive_scalar} depends only on the
surface-curl/Faraday content of the shell tangential electric field. Pure
curl-free shell electric fields satisfy
\[
c = 0
\quad\Longrightarrow\quad
f_{\rm ind}=0.
\]
So the condensed operator is naturally an \emph{inductive} forcing object of
the reduced tangential closure, not an instantaneous shell-current continuity
object. This is why it remains nontrivial exactly where the shell-current
candidate can be too permissive.

\paragraph{Reduced-model status.}
The exact projected Maxwell driver is
\[
f_t^{\rm drv}
=
\mathbf B_{0S}\cdot\bigl[-\nabla\times\nabla\times \mathbf E^{\rm drv}\bigr]_S.
\]
The live operator \eqref{eq:condensed_forcing_toroidal_vector_form} is a
reduction of this exact driver: it keeps the part built from the radial curl
scalar of the shell tangential field and drops the remaining tangential
curlcurl content. This is why it is best described as the inductive forcing
object of the current projected tangential closure, rather than as the exact
radial-shell trace.

\paragraph{A frozen-conductance incremental law.}
There is, however, a clean first-principles law for the forcing-side scalar.
Interpret the forcing input as an \emph{increment} in the shell tangential
electric field over the instantaneous connector substep, and assume the shell
conductance tensor is frozen during that substep. Then
\[
\delta \mathbf K_S = \Sigma\,\delta \mathbf E_S,
\]
or, in the discrete coefficient formulation, by inverting the live
\(\mathbf K_S \mapsto \mathbf E_S\) operator on the active shell space.
Thin-sheet current continuity then gives
\[
\delta j_r^+
=
-\,\nabla_S\cdot \delta \mathbf K_S,
\]
so the forcing-side radial-shell scalar increment is
\begin{equation}
\delta f_{\rm known}
=
\mu_0\,\delta j_{r,\rm known}^+
=
-\,\frac{\mu_0}{R_I}\,\mathrm{div}_\Omega(\delta \mathbf K_{S,\rm known}).
\label{eq:forcing_shell_continuity_incremental}
\end{equation}
This gives an explicit shell-law-based forcing operator
\[
\delta\mathbf E_{S,\rm known}
\;\to\;
\delta\mathbf K_{S,\rm known}
\;\to\;
\delta f_{\rm known}
\]
without routing through the shadow tangential toroidal solve.

\paragraph{Current-first interpretation.}
The most coherent reading of \eqref{eq:forcing_shell_continuity_incremental}
is therefore not ``electric-first'' but ``current-first'':
\[
\delta\mathbf K_{S,\rm known}
\;\to\;
\delta j_{r,\rm known}^+
\]
is the primitive shell law, and the frozen-conductance constitutive relation
\[
\delta\mathbf E_{S,\rm known}\mapsto \delta\mathbf K_{S,\rm known}
\]
is the adapter that makes this branch usable with the current runtime forcing
state. This matches the feedback-side shell-current viewpoint more directly
than treating \(\delta\mathbf E_S\mapsto \delta j_r^+\) as the primitive law.

\paragraph{If \(\partial_t\Sigma=0\).}
Under the stronger assumption that the conductance tensor is static in time,
the same constitutive step can be written in differential form as
\[
\partial_t \mathbf K_S = \Sigma\,\partial_t \mathbf E_S.
\]
The present note still writes the forcing-side branch in \(\delta\)-form,
because the runtime input is currently interpreted as a connector-substep
increment rather than as an explicitly exposed \(\partial_t \mathbf E_S\)
field. But the shell-current source law itself is the same.

\paragraph{Most intuitive reading of the incremental forcing law.}
This is the same pillbox argument, but applied to the \emph{known forcing}
branch alone. In that branch there is no independent prescribed lower-side FAC
source. So any lateral sheet-current increment created by the imposed shell
electric increment must close by leaving the shell upward:
\[
\delta j_{r,\rm known}^+
=
-\,\nabla_S\!\cdot \delta\mathbf K_{S,\rm known}.
\]
In words,
\[
\text{upward forcing-driven current}
=
\text{minus the lateral current convergence created in the sheet}.
\]
This is why \eqref{eq:forcing_shell_continuity_incremental} is the most
natural first-principles law if the input data are genuinely interpreted as
\(\delta\mathbf E_S\).

\paragraph{Historical instantaneous reading.}
The same matrix can be written formally by recovering the shell current from the
instantaneous shell constitutive law,
\[
\mathbf K_{S,\rm known} = \Sigma\,\mathbf E_{S,\rm known},
\]
but that reading is only a reduced approximation because the radial-shell
equation is an equation for a time derivative.

\paragraph{Status of that incremental law.}
The current code now exposes \eqref{eq:forcing_shell_continuity_incremental} as
an explicit alternative forcing model. Under the frozen-conductance
incremental interpretation it is first-principles. But on representative
pure-spectral full-induction states it does \emph{not} coincide with the
previously implemented condensed forcing operator. So this forcing-side law
should be regarded as a serious first-principles alternative, not yet as a
proven replacement for the condensed PFAC/tangential forcing branch.

\paragraph{Where the real fragility lies.}
The main uncertainty in \eqref{eq:forcing_shell_continuity_incremental} is not
the use of thin-sheet current continuity itself. Within the adopted shell model
that part is the same robust pillbox balance used on the feedback side. The
more fragile ingredients are:
\begin{enumerate}
\item interpreting the supplied forcing data as the actual shell-electric
increment \(\delta\mathbf E_S\) over the connector substep;
\item freezing the shell conductance tensor during that substep; and
\item neglecting any stored surface-charge contribution in that forcing-side
incremental balance.
\end{enumerate}
So if one is uneasy about this branch, the first place to look is the forcing
semantics, not the continuity identity itself.

\paragraph{Exact q-trace representation of the incremental law.}
In the newer forcing-side architecture, this same shell-current law is now also
represented through the exact \(q\)-trace interface:
\[
\delta \mathbf E_{S,\rm known}
\mapsto
\delta \mathbf K_{S,\rm known}
\mapsto
\delta j_{r,\rm known}^+
\mapsto
\delta\psi_{\rm known}^+
\mapsto
q_{\rm known}=R_I\,\delta\psi_{\rm known}^+.
\]
So even the explicit frozen-conductance branch now feeds the radial-shell
closure through the exact shell quantity \(q_{\rm known}=\partial_rU-E_r\).
The semantic difference between the built-in forcing modes is therefore not
\emph{whether} they use the exact shell inversion, but \emph{which closure law}
is used to build the source \(\partial_t j_{r,\rm known}^+\).

\paragraph{Is that mismatch a bug?}
Not by itself. The condensed forcing operator and the frozen-conductance
incremental law are \emph{different closure statements}. The condensed operator
keeps only the inductive radial-curl/toroidal content of the shell tangential
electric field, whereas \eqref{eq:forcing_shell_continuity_incremental}
interprets the input as an actual shell-electric \emph{increment} and then
propagates it through instantaneous sheet-current continuity. If those two
operators fail to match, that is only a contradiction if one claims that the
same runtime quantity \(\mathbf E_{S,\rm known}\) is simultaneously:
\begin{itemize}
\item the inductive forcing data of the reduced projected-tangential driver; and
\item the literal shell-electric increment \(\delta \mathbf E_S\) of the
frozen-conductance connector substep.
\end{itemize}
The present code does \emph{not} make that identification. So the observed
mismatch should currently be read as a difference in closure semantics, not as
an automatic algebraic bug.

There is also an important interpretive caution. The shell constitutive law
\[
\mathbf K_S = \Sigma \mathbf E_S
\]
relates the instantaneous shell current to the instantaneous shell electric
field. The radial-shell equation, however, is an equation for
\[
\partial_t j_r.
\]
So the feedback-side continuity law is stronger: there the runtime genuinely
provides
\[
\partial_t\psi \to \partial_t\mathbf K_S,
\]
and continuity may be differentiated consistently. On the forcing side, the
known data are \(\mathbf E_{S,\rm known}\), not \(\partial_t\mathbf E_S\) or
\(\partial_t\mathbf K_S\). Therefore the explicit candidate
\eqref{eq:forcing_shell_continuity_incremental} should be viewed as an
incremental/reduced-model forcing law, not as a settled exact identity of the
same status as the feedback-side continuity connector.

If surface-charge storage or other non-quasi-static sheet effects were thought
to matter over the connector substep, then this forcing-side balance would need
to be enlarged before the frozen-conductance incremental law could be treated
as complete.

This caution is not merely formal. Direct operator diagnostics in the current
runtime show that the explicit shell-current candidate responds strongly to the
curl-free shell-electric component as well as to the divergence-free one in a
simple wind-driven test case, whereas the older condensed induction forcing
operator is nearly null on both components there. So the explicit candidate is
too permissive in exactly the way one would expect from an instantaneous
\(\mathbf K_S=\Sigma\mathbf E_S\) current law: it is sensitive to shell
electric structure that need not represent the inductive
\(\partial_t j_r\)-driving content of the forcing.

A more conservative time-derivative candidate can also be written. Since the
runtime already exposes
\[
E_{\mathrm{df}} \mapsto \partial_t m_{\mathrm{ind}},
\]
one may try to continue that through the induced poloidal sheet-current map,
\[
\partial_t m_{\mathrm{ind}}
\mapsto
\partial_t \mathbf K_{S,\mathrm{ind}}
\mapsto
-\,\nabla_S\!\cdot \partial_t \mathbf K_{S,\mathrm{ind}}.
\]
But in the present operator chain this candidate collapses to zero: the
divergence of the induced-polodial sheet-current branch vanishes in the active
coefficient formulation. So that route does not supply the missing forcing-side
scalar either.

Taken together, these diagnostics indicate that the forcing-side scalar should
still be interpreted as an \emph{inductive} object built from the Faraday/curl
content of the shell electric response, not from the instantaneous shell
current alone. In practice, this leaves the existing condensed
curlcurl/Faraday-based forcing operator as the only nontrivial forcing-side
object currently supported by the runtime. Operationally, the built-in default
forcing scalar is therefore the composition
\eqref{eq:built_in_condensed_forcing_chain}, not an independent shell-current
continuity law.

\paragraph{No double counting.}
The split between forcing and feedback is intentional. In the runtime
\texttt{radial\_shell} assembly:
\begin{itemize}
\item the forcing-side scalar response comes only from the condensed nonlocal
PFAC response route;
\item the induced \texttt{dt\_alpha} feedback comes only from the explicit
thin-shell current-continuity connector route.
\end{itemize}
So the runtime is not summing two overlapping approximations of the same
radial-shell scalar. It is composing two distinct pieces of the same linear
response to a trial toroidal update.

\paragraph{Raw connector versus continuity connector.}
The raw explicit $R_M$ toroidal boundary connector by itself is \emph{not} the
final radial-shell scalar. Direct operator comparison on representative
full-induction spectral states shows a material mismatch between
\[
\partial_t\alpha \to \partial_t\psi^+_{\rm raw}
\to -R_I\Delta_S(\partial_t\psi^+_{\rm raw})
\]
and
\[
\partial_t\alpha \to \text{self-consistent shell-feedback scalar}.
\]
What survives that comparison is the thin-shell current-continuity law
\eqref{eq:thin_shell_current_continuity_connector}: when the explicit shell
current operator
\[
\partial_t\psi \to \partial_t\mathbf K_S
\]
is used, the resulting continuity connector reproduces the previously matched
corrected connector to numerical precision on representative spectral states.
So the present \texttt{radial\_shell} runtime no longer needs the raw
\(R_M\)-toroidal branch as the feedback closure; the feedback closure is the
explicit thin-shell current-continuity law, understood within the adopted
instantaneous thin-sheet model with no stored surface charge over the connector
substep.

\section{What part of this is exact, and what part is modeled}

\subsection*{Exact within the adopted shell model}

The following parts are exact consequences of the shell identities, together
with the adopted thin-sheet shell model where indicated:
\begin{enumerate}
\item $T_{\alpha\to j_r}$ from $j_r = B_{0r}\alpha$;
\item $T_{j_r\to\psi}$ from
$j_r = -(\mu_0 R_I)^{-1}\Delta_\Omega \psi$;
\item the exact scalar equivalence
\eqref{eq:missing_scalar_equivalence};
\item the explicit thin-shell current-continuity connector
\eqref{eq:thin_shell_current_continuity_connector} used to construct
$\partial_t\psi^+$ or $\partial_t j_r^+$ from the shell source and the
explicit \(\partial_t\psi\to\partial_t\mathbf K_S\) map, within the adopted
thin-sheet/no-stored-charge connector substep model.
\end{enumerate}

\subsection*{Still modeled rather than implied by Maxwell alone}

The present built-in radial-shell runtime is stronger than the earlier
equivalent tangential adapter, but it is still an explicit model, not the full
gap-region PDE solve. The remaining model assumptions are:
\begin{enumerate}
\item the forcing-side radial-shell scalar is taken from the reduced
tangential closure chain \eqref{eq:built_in_condensed_forcing_chain}, rather
than from the exact radial trace \(Q_I-\Delta_\Omega E_{r,I}\);
\item within that reduced tangential forcing chain, the live driver keeps only
the toroidal contribution
\eqref{eq:condensed_forcing_toroidal_vector_form} of the full projected
tangential Maxwell driver;
\item the chosen open/closed $R_M$ policy for the auxiliary PFAC response and
other upper-region response operators;
\item optional shell-electric trace models such as harmonic continuation of
\(U\) and ideal-gap recovery of \(E_r\) remain available explicitly for
diagnostics, but they are no longer part of the built-in default path.
\end{enumerate}

The last item belongs only to the optional explicit trace-based route; it is
not part of the built-in default \texttt{radial\_shell} path anymore. The
first three items are still part of the default reduced PFAC/RM shell model.
So the runtime now has an \emph{explicit} PFAC/RM radial-shell response model
without harmonic \(U\) in the default path, but it is still not the same thing
as solving the full upper-region boundary value problem in the gap directly.

\paragraph{What ``full'' means here.}
Relative to the adopted model assumptions already used elsewhere in the code
(PFAC/radial current closure, chosen open/closed $R_M$ semantics where those
operators appear, and the thin-shell current-continuity connector), the runtime now
has a \emph{closed linear radial-shell response model} for that adopted
PFAC/RM semantics. An optional trace-based variant with harmonic continuation of the
poloidal electric scalar and ideal-gap recovery of $E_r$ also exists, but it
is no longer the default built-in interpretation.
What it does \emph{not} claim
is an assumption-free direct solve of the gap-region Maxwell problem with no
additional constitutive choices.

\section{Comparison with the older tangential closures}

The code base now has one supported toroidal closure:
\begin{enumerate}
\item \textbf{\path{radial_shell}:} explicit PFAC/RM shell-gap closure in
scalar shell variables, built around the connector \(\chi\), the shell
magnetic-return operator, and the condensed shell response.
\end{enumerate}

Older tangential closures remain relevant only as internal historical
background and low-level benchmark machinery; they are no longer supported
runtime model choices.

The projected tangential analysis above remains useful because:
\begin{itemize}
\item it still explains where the $|\mathbf B_{0S}|^2$ operator comes from;
\item it still shows that the projected scalar tangential closure is weaker
than the full tangential vector equation;
\item and it still leaves the sign of the assembled projected feedback block as
an open issue for the tangential path.
\end{itemize}

\section{Recommended interpretation for the code base}

The formal documentation should now say:
\begin{enumerate}
\item The canonical scalar first-principles closure remains the radial-shell
equation \eqref{eq:closed_radial_shell_response_problem}.

\item The runtime \path{radial_shell} mode is now an explicit PFAC/RM-based
radial-shell response model. It should be documented as such.

\item The maps
$\eta \to \partial_t j_r \to \partial_t \psi \to \partial_t\psi^+$
are tied to exact shell identities plus the explicit thin-shell
current-continuity connector in the built-in default path.

\item The built-in \path{radial_shell} runtime should document its
forcing-side choice explicitly: the supported built-in realization is the
incremental shell-current law
\eqref{eq:forcing_shell_continuity_incremental}. Older condensed or
homogeneous-outer forcing constructions can still be written down explicitly,
but they are no longer part of the supported built-in mode surface.

\item The explicit frozen-conductance incremental shell-current law should be
documented as an alternative forcing closure. Its mismatch with the condensed
forcing operator should \emph{not} be described as a bug by default, because
the two operators act on differently interpreted forcing data.

\item The optional forcing-side shell electric trace route should be documented
as an explicit upper-region model:
\[
[\Phi, W] \mapsto \bigl(\partial_r U|_{R_I^+}, E_{r,I}\bigr),
\]
with its assumptions stated directly.

\item The tangential projected and tangential full modes remain alternative
closures. They should not be conflated with the radial-shell formulation.

\item The sign of the assembled projected tangential feedback block remains a
separate issue for the tangential path; the present note does not resolve it.

\item The least-squares solve, low-latitude coupling constraints, and
regularized $j_r\to\alpha$ inversion remain numerical/model layers on top of
the exact shell identities.
\end{enumerate}

\end{document}
